% ============================================================
%  SECCIÓN 3 — Descripción y Exploración de Datos
%  Sistema de Alerta Temprana (SAT) — IIND-2104
%  Universidad de los Andes
%  Autores: Nicolás Torres Pulido, Isabella Delgadillo Calero
%  Director: Juan Fernando Pérez Bernal
%  Citación: IEEE
%
%  NOTA DE INTEGRACIÓN: Para incluir en el documento principal,
%  usar % ============================================================
%  SECCIÓN 3 — Descripción y Exploración de Datos
%  Sistema de Alerta Temprana (SAT) — IIND-2104
%  Universidad de los Andes
%  Autores: Nicolás Torres Pulido, Isabella Delgadillo Calero
%  Director: Juan Fernando Pérez Bernal
%  Citación: IEEE
%
%  NOTA DE INTEGRACIÓN: Para incluir en el documento principal,
%  usar % ============================================================
%  SECCIÓN 3 — Descripción y Exploración de Datos
%  Sistema de Alerta Temprana (SAT) — IIND-2104
%  Universidad de los Andes
%  Autores: Nicolás Torres Pulido, Isabella Delgadillo Calero
%  Director: Juan Fernando Pérez Bernal
%  Citación: IEEE
%
%  NOTA DE INTEGRACIÓN: Para incluir en el documento principal,
%  usar % ============================================================
%  SECCIÓN 3 — Descripción y Exploración de Datos
%  Sistema de Alerta Temprana (SAT) — IIND-2104
%  Universidad de los Andes
%  Autores: Nicolás Torres Pulido, Isabella Delgadillo Calero
%  Director: Juan Fernando Pérez Bernal
%  Citación: IEEE
%
%  NOTA DE INTEGRACIÓN: Para incluir en el documento principal,
%  usar \input{seccion3_exploracion_datos} y eliminar el
%  preámbulo y los comandos \begin{document}/\end{document}.
%  Cambiar [english] a [spanish] con texlive-lang-spanish.
%  Las figuras se cargan desde la carpeta figures/ relativa.
% ============================================================

\documentclass[12pt, letterpaper]{article}

\usepackage[utf8]{inputenc}
\usepackage[T1]{fontenc}
\usepackage[english]{babel}
\usepackage{amsmath, amssymb}
\usepackage{geometry}
\usepackage{setspace}
\usepackage{booktabs}
\usepackage{multirow}
\usepackage{array}
\usepackage{graphicx}
\usepackage{subcaption}
\usepackage{xcolor}
\usepackage{cite}
\usepackage[hidelinks]{hyperref}
\usepackage{url}
\usepackage{rotating}

\geometry{top=2.5cm, bottom=2.5cm, left=3cm, right=2.5cm}
\onehalfspacing
\renewcommand{\refname}{Referencias}

\graphicspath{{figures/}}

% ============================================================
\begin{document}
% ============================================================

\section{Descripción y Exploración de Datos}
\label{sec:exploracion}

Este capítulo caracteriza el conjunto de datos empleado en el SAT,
documenta las decisiones de preprocesamiento adoptadas y presenta
los hallazgos del análisis exploratorio que fundamentan el diseño
del sistema de modelado. La presentación se organiza en seis
bloques: (i)~fuente y composición del dataset;
(ii)~definición y justificación de los subconjuntos de trabajo;
(iii)~calidad y estructura de los datos;
(iv)~análisis estadístico de las notas académicas;
(v)~análisis de la actividad en plataforma; y
(vi)~análisis de correlaciones y síntesis de decisiones de diseño.

% ------------------------------------------------------------
\subsection{Fuente y Composición del Conjunto de Datos}
\label{subsec:fuente}
% ------------------------------------------------------------

Los datos provienen del curso IIND-2104 \textit{Modelos
Probabilísticos} de la Universidad de los Andes, dictado entre
los semestres 202320 y 202520 —cinco cohortes consecutivas
cubiertas entre 2023 y 2025. La información se integra a partir
de dos fuentes complementarias:

\begin{itemize}
  \item \textbf{Sistema de gestión académica institucional:}
    calificaciones oficiales de los tres parciales y de las
    evaluaciones continuas (actividades modulares, quizzes,
    talleres y proyectos, según el formato de cada semestre).

  \item \textbf{Plataforma Bloque Neón:} sistema de gestión del
    aprendizaje (LMS) que registra, por sesión y por módulo
    de contenido, el tiempo de estudio, los temas consultados,
    los módulos visitados y la frecuencia de acceso de cada
    estudiante.
\end{itemize}

El dataset crudo consolidado cuenta con \textbf{764~registros y
77~variables}, con cuatro posibles resultados académicos finales:
\texttt{APROBÓ} (413 estudiantes, 54,1\,\%),
\texttt{REPROBÓ} (141, 18,5\,\%),
\texttt{RETIRADO} (164, 21,5\,\%) y
\texttt{SIN~NOTA} (46, 6,0\,\%).

La Figura~\ref{fig:composicion} ilustra cómo se distribuyen
estos resultados por semestre en términos absolutos y
porcentuales, y la Tabla~\ref{tab:composicion_semestral} detalla
las frecuencias exactas.

\begin{figure}[htbp]
  \centering
  \includegraphics[width=\textwidth]{01_composicion_dataset.png}
  \caption{Composición del dataset por semestre y resultado
           académico. \textbf{Izquierda:} distribución absoluta.
           \textbf{Derecha:} distribución porcentual con
           etiquetas de proporción. El sombreado naranja
           corresponde a la era pedagógica antigua (202320--202410)
           y el verde a la nueva (202420--202520).}
  \label{fig:composicion}
\end{figure}

\begin{table}[htbp]
  \centering
  \caption{Composición por semestre, era pedagógica y resultado
           académico final.}
  \label{tab:composicion_semestral}
  \small
  \begin{tabular}{lccrrrrrrr}
    \toprule
    \textbf{Semestre} & \textbf{Era} & \textbf{Total}
      & \multicolumn{2}{c}{\textbf{APROBÓ}}
      & \multicolumn{2}{c}{\textbf{REPROBÓ}}
      & \multicolumn{2}{c}{\textbf{RETIRADO}}
      & \textbf{SIN} \\
    \cmidrule(lr){4-5}\cmidrule(lr){6-7}\cmidrule(lr){8-9}
    & & & $n$ & \% & $n$ & \% & $n$ & \% & \textbf{NOTA} \\
    \midrule
    202320 & Antigua & 170 & 116 & 68,2 &  27 & 15,9 &  27 & 15,9 & 0  \\
    202410 & Antigua & 150 &  68 & 45,3 &  14 &  9,3 &  30 & 20,0 & 38 \\
    202420 & Nueva   & 164 &  90 & 54,9 &  33 & 20,1 &  40 & 24,4 & 1  \\
    202510 & Nueva   & 140 &  70 & 50,0 &  47 & 33,6 &  21 & 15,0 & 2  \\
    202520 & Nueva   & 140 &  69 & 49,3 &  20 & 14,3 &  46 & 32,9 & 5  \\
    \midrule
    \textbf{Total} & — & \textbf{764}
      & \textbf{413} & \textbf{54,1}
      & \textbf{141} & \textbf{18,5}
      & \textbf{164} & \textbf{21,5}
      & \textbf{46} \\
    \bottomrule
  \end{tabular}
\end{table}

Dos fenómenos merecen atención especial. Primero, la tasa de
reprobación en 202510 (33,6\,\%) es la más alta del período y
más del doble de la de 202410 (9,3\,\%), lo que sugiere
variabilidad intercohorte que el modelo debe generalizar. Segundo,
el semestre 202520 registra la mayor tasa de retiro (32,9\,\%),
fenómeno atípico que podría obedecer a cambios en la política
institucional de permanencia o a condiciones externas no
capturadas en los datos.

% ............................................................
\subsubsection{Dos eras pedagógicas}
\label{sssec:eras}
% ............................................................

Un rasgo estructural del dataset es la coexistencia de dos
\textbf{formatos de evaluación continua}:

\begin{itemize}
  \item \textbf{Era Antigua} (202320 y 202410; $n = 320$ estudiantes
    activos): las evaluaciones continuas consisten en talleres
    individuales y proyectos por fases grupales.

  \item \textbf{Era Nueva} (202420, 202510 y 202520; $n = 444$
    estudiantes activos): las evaluaciones continuas son
    actividades modulares (AMs) y quizzes individuales.
\end{itemize}

Dentro de la era nueva existe una distinción adicional: en
202420 las AMs se realizaron \textbf{en grupo}, mientras que
a partir de 202510 pasaron a ser \textbf{individuales}. Este
cambio produce diferencias estadísticamente significativas en la
distribución de AM-3 ($U$-test, $p = 1{,}39 \times 10^{-9}$,
$d_{\text{Cohen}} = 0{,}40$, AUC$_{\text{grupo}} = 0{,}369$
vs.\ AUC$_{\text{ind.}} = 0{,}246$), lo que justifica incluir la
variable \texttt{modalidad\_am} como control en los modelos
entrenados con datos del formato nuevo.

% ------------------------------------------------------------
\subsection{Definición y Justificación de los Subconjuntos
            de Trabajo}
\label{subsec:subsets}
% ------------------------------------------------------------

El dataset bruto requiere dos depuraciones sucesivas antes de
ser apto para el modelado de clasificación.

% ............................................................
\subsubsection{Exclusión de estudiantes retirados y sin
               calificación}
\label{sssec:excl_retirados}
% ............................................................

Los 164 estudiantes con resultado \texttt{RETIRADO} y los
46 con \texttt{SIN~NOTA} se excluyen por tres razones:

\begin{enumerate}
  \item \textbf{Ausencia de etiqueta objetivo observable.}
    El SAT modela $P(\text{REPROBÓ} \mid \mathbf{x}_i)$, la
    probabilidad de que un estudiante \emph{finalice} el semestre
    reprobando. Los estudiantes retirados abandonaron el curso
    antes de producirse el resultado; su etiqueta final es
    desconocida y no puede inferirse sin supuestos fuertes.

  \item \textbf{Sesgo de retiro endógeno.}
    La decisión de retirarse está correlacionada con el
    desempeño anticipado: los estudiantes que perciben que van
    a reprobar se retiran en mayor proporción \cite{Tinto1993}.
    Incluirlos en la clase \texttt{REPROBÓ} introduciría sesgo
    de selección en las variables de engagement: un estudiante
    retirado en la semana~4 tiene engagement mínimo no por falta
    de habilidad sino por abandono temprano, contaminando el
    predictor más utilizado.

  \item \textbf{Heterogeneidad semántica.}
    La predicción del retiro como fenómeno independiente
    requeriría modelos multi-clase y supuestos adicionales
    sobre su equivalencia con la reprobación, lo cual está
    fuera del alcance del presente trabajo.
\end{enumerate}

El subconjunto resultante es \textbf{\texttt{df\_activos}}
($N = 554$; APROBÓ: 413, REPROBÓ: 141).

% ............................................................
\subsubsection{Exclusión adicional de estudiantes sin actividad
               en plataforma}
\label{sssec:excl_engagement}
% ............................................................

Del subconjunto \texttt{df\_activos}, 52 estudiantes (9,4\,\%)
presentan \textbf{engagement total nulo} en Bloque Neón durante
todo el semestre (i.e., \texttt{engagement\_hasta\_p1} $= 0$).
Estos casos se excluyen del subconjunto de modelado por las
razones siguientes:

\begin{enumerate}
  \item \textbf{Ambigüedad del cero.}
    Existen dos interpretaciones mutuamente excluyentes:
    (a)~el estudiante genuinamente no utilizó la plataforma,
    o (b)~ocurrió un error en el pipeline de integración de
    datos. Ambas hipótesis son plausibles: 39 de los 52 casos
    (75\,\%) se concentran en el semestre 202320, el primero
    en incorporar Bloque Neón, cuando el sistema estaba aún
    en fase de implementación. No es posible distinguir entre
    causas con la información disponible.

  \item \textbf{Riesgo de imputación espuria.}
    Imputar cero como valor real de engagement asumiría que
    ``no usar la plataforma'' es una señal homogénea entre
    semestres, cuando los patrones de no-uso varían
    considerablemente entre cohortes.

  \item \textbf{Impacto estadístico marginal confirmado.}
    Se realizó un análisis de sensibilidad: los modelos
    entrenados sobre \texttt{df\_activos} completo ($N = 554$)
    y sobre el subconjunto filtrado ($N = 502$) arrojan
    diferencias en F1 menores a $\Delta F_1 < 0{,}02$,
    confirmando que la exclusión no altera las conclusiones.
\end{enumerate}

El subconjunto de modelado final es
\textbf{\texttt{df\_activos\_engagement}}
($N = 502$; APROBÓ:~373, REPROBÓ:~129).
La Tabla~\ref{tab:subsets} resume los tres subconjuntos
empleados en el análisis.

\begin{table}[htbp]
  \centering
  \caption{Subconjuntos de trabajo. \texttt{df\_activos\_engagement}
           es el conjunto empleado para el entrenamiento y
           evaluación de todos los modelos del SAT.}
  \label{tab:subsets}
  \begin{tabular}{lrrrl}
    \toprule
    \textbf{Subconjunto} & \textbf{$N$}
      & \textbf{APROBÓ} & \textbf{REPROBÓ}
      & \textbf{Uso principal} \\
    \midrule
    \texttt{df\_all}
      & 764 & 413 & 141
      & Tasas globales, análisis de retiro \\
    \texttt{df\_activos}
      & 554 & 413 & 141
      & Análisis de notas y distribuciones \\
    \texttt{df\_activos\_engagement}
      & \textbf{502} & \textbf{373} & \textbf{129}
      & \textbf{Entrenamiento y evaluación de modelos} \\
    \bottomrule
  \end{tabular}
\end{table}

% ------------------------------------------------------------
\subsection{Calidad y Estructura de los Datos}
\label{subsec:calidad}
% ------------------------------------------------------------

% ............................................................
\subsubsection{Nulos estructurales versus problemáticos}
\label{sssec:nulos}
% ............................................................

El dataset presenta un volumen considerable de valores faltantes,
en su mayoría atribuibles al cambio de formato pedagógico.
Se distinguen dos categorías.

\paragraph{Nulos estructurales (esperados).}
Son variables que solo existen en una era y, por tanto, están
ausentes de forma legítima en la otra:
\texttt{quiz\_4/5/6} (77,8\,\% de nulos, solo en algunos
semestres del formato nuevo),
\texttt{taller\_1/2/3} (59,4\,\%, exclusivo de la era antigua),
\texttt{fase\_1/2} (59,4\,\%, proyectos grupales de la era antigua),
y \texttt{am\_4/5/6} (56,7\,\%, AMs opcionales no presentes en
todos los semestres). Estos nulos no representan problemas de
calidad sino de diseño del instrumento de evaluación.

\paragraph{Nulos problemáticos (casos especiales).}
Son faltantes en variables que deberían estar presentes para
todos los estudiantes activos:
\begin{itemize}
  \item \textbf{Parcial~1 $= 0$ ($n = 1$):} un único estudiante
    registra nota cero. No es posible distinguir ausencia de nota
    real. Se retiene con variable indicadora
    \texttt{no\_presento\_p1}.
  \item \textbf{Parcial~3 $= 0$ ($n = 15$; 13 reprobaron):}
    concentrado en estudiantes encaminados a reprobar, consistente
    con ausentismo por desmotivación al final del semestre.
    Se retienen con su valor real.
  \item \textbf{Notas superiores a 5{,}0:} presentes por
    bonificaciones otorgadas por docentes. Son valores legítimos
    y se retienen sin transformación.
\end{itemize}

% ............................................................
\subsubsection{Variables de plataforma sin poder discriminativo}
\label{sssec:excl_vars}
% ............................................................

La plataforma Bloque Neón registra tanto métricas de
\textit{estudio de contenidos} como métricas de
\textit{actividades evaluativas en línea}. El análisis estadístico
univariante sobre estas últimas —\texttt{actividades\_n\_realizadas},
\texttt{actividades\_tiempo\_total} y
\texttt{actividades\_visitas\_total}— revela ausencia de
poder discriminativo entre aprobados y reprobados
(Tabla~\ref{tab:excl_vars}), lo que se ilustra también en la
Figura~\ref{fig:vars_excluidas}.

\begin{table}[htbp]
  \centering
  \caption{Variables de actividades evaluativas de la plataforma.
           AUC~$\approx 0{,}5$ y $p > 0{,}05$ indican que no
           discriminan mejor que la clasificación aleatoria.}
  \label{tab:excl_vars}
  \begin{tabular}{lrrrr}
    \toprule
    \textbf{Variable}
      & \textbf{Med.\ APROBÓ}
      & \textbf{Med.\ REPROBÓ}
      & \textbf{AUC}
      & \textbf{$p$-valor (MW)} \\
    \midrule
    \texttt{actividades\_n\_realizadas}   & 1,00 & 1,00 & 0,473 & 0,830 \\
    \texttt{actividades\_tiempo\_total}   & 1,15 & 1,10 & 0,498 & 0,532 \\
    \texttt{actividades\_visitas\_total}  & 9,00 & 11,00 & 0,465 & 0,886 \\
    \bottomrule
  \end{tabular}
\end{table}

\begin{figure}[htbp]
  \centering
  \includegraphics[width=\textwidth]{11_variables_excluidas.png}
  \caption{Distribuciones de las tres variables de actividades
           evaluativas para aprobados (azul) y reprobados (rojo).
           Los valores de $p > 0{,}05$ y AUC~$< 0{,}5$
           confirman la ausencia de poder discriminativo,
           justificando su exclusión del conjunto de predictores.}
  \label{fig:vars_excluidas}
\end{figure}

Nótese que, para \texttt{actividades\_visitas\_total}, la mediana
de los reprobados (11) supera la de los aprobados (9), con
AUC~$= 0{,}465$ —peor que la clasificación aleatoria—, lo que
sugiere que los estudiantes con mayor número de visitas a
actividades no necesariamente obtienen mejores resultados. Las
tres variables son excluidas del conjunto de predictores para
evitar introducir ruido en los modelos.

% ------------------------------------------------------------
\subsection{Análisis Estadístico de las Notas Académicas}
\label{subsec:notas}
% ------------------------------------------------------------

El análisis de notas se realiza sobre \texttt{df\_activos}
($N = 554$) para preservar la distribución real sin los
filtros de engagement.

% ............................................................
\subsubsection{Distribución de la nota final}
\label{sssec:dist_nota}
% ............................................................

La nota final (escala 0--5) presenta una distribución bimodal
con dos concentraciones claramente separadas: una en torno a
2,6 (zona de reprobación) y otra en torno a 3,7 (zona de
aprobación cómoda), como se observa en la
Figura~\ref{fig:dist_nota}.

\begin{figure}[htbp]
  \centering
  \includegraphics[width=\textwidth]{02_distribucion_nota_final.png}
  \caption{Distribución de la nota final.
           \textbf{Izquierda:} histogramas superpuestos con
           curva de densidad KDE por resultado académico.
           La línea discontinua vertical indica el umbral de
           aprobación (3{,}0).
           \textbf{Derecha:} diagramas de caja por semestre y
           resultado; el sombreado distingue las eras pedagógicas.}
  \label{fig:dist_nota}
\end{figure}

Las estadísticas descriptivas detalladas se presentan en la
Tabla~\ref{tab:desc_nota}. La prueba de Mann-Whitney confirma
la separación altamente significativa entre grupos
($U = 58\,233{,}0$, $p = 9{,}97 \times 10^{-71}$,
$d_{\text{Cohen}} = 2{,}41$). Esta magnitud del efecto implica
que la distribución de notas de los aprobados y la de los
reprobados se solapan apenas en una pequeña región central
(alrededor de $[2{,}8;\;3{,}2]$), que es precisamente la zona
donde los modelos de clasificación concentran su incertidumbre.

\begin{table}[htbp]
  \centering
  \caption{Estadísticas descriptivas de la nota final por resultado
           y era pedagógica. $\bar{x}$: media; $\tilde{x}$: mediana;
           $s$: desviación estándar.}
  \label{tab:desc_nota}
  \begin{tabular}{lrrrr}
    \toprule
    \textbf{Grupo} & \textbf{$N$}
      & \textbf{$\bar{x}$} & \textbf{$\tilde{x}$} & \textbf{$s$} \\
    \midrule
    APROBÓ (global)                & 413 & 3,724 & 3,620 & 0,527 \\
    REPROBÓ (global)               & 141 & 2,418 & 2,632 & 0,583 \\
    \midrule
    Era Antigua — APROBÓ           & 184 & 3,777 & 3,680 & 0,504 \\
    Era Antigua — REPROBÓ          &  41 & 2,446 & 2,720 & 0,671 \\
    Era Nueva   — APROBÓ           & 229 & 3,682 & 3,530 & 0,541 \\
    Era Nueva   — REPROBÓ          & 100 & 2,407 & 2,605 & 0,545 \\
    \bottomrule
  \end{tabular}
\end{table}

% ............................................................
\subsubsection{Poder predictivo de los parciales:
               análisis por era pedagógica}
\label{sssec:parciales_era}
% ............................................................

La Figura~\ref{fig:parciales_era} compara la distribución de
las tres calificaciones de parciales entre aprobados y reprobados,
estratificada por era pedagógica.

\begin{figure}[htbp]
  \centering
  \includegraphics[width=\textwidth]{05_parciales_era_resultado.png}
  \caption{Distribuciones de Parcial~1, Parcial~2 y Parcial~3
           por resultado académico, estratificadas por era
           pedagógica. Cada subgráfico incluye el $p$-valor
           de la prueba Mann-Whitney, el AUC univariante y el
           coeficiente de rango biserial~$r_b$.}
  \label{fig:parciales_era}
\end{figure}

La Tabla~\ref{tab:auc_era} sintetiza los AUC univariantes por
variable y era. El hallazgo más relevante es la \textbf{reducción
del poder predictivo de P1 en la era nueva}
($\Delta\text{AUC} = -0{,}087$). Esta caída puede atribuirse
al mayor peso relativo de las actividades continuas (quizzes y
AMs) en el nuevo currículo, que diluye la señal del primer
examen. Contrariamente, el Parcial~2 mantiene una capacidad
predictiva alta y estable en ambas eras
($\Delta\text{AUC} = -0{,}025$), lo que refuerza el valor
del segundo checkpoint (CP2) del SAT.

\begin{table}[htbp]
  \centering
  \caption{AUC univariante de las variables académicas por era
           pedagógica. $\Delta$AUC indica la variación entre eras.
           Diferencias $> 5\,\%$ se consideran sustanciales.}
  \label{tab:auc_era}
  \begin{tabular}{lrrrc}
    \toprule
    \textbf{Variable}
      & \textbf{AUC Era Antigua}
      & \textbf{AUC Era Nueva}
      & \textbf{$\Delta$AUC}
      & \textbf{Sustancial} \\
    \midrule
    Parcial~1      & 0,869 & 0,781 & $-0{,}087$ & Sí \\
    Parcial~2      & 0,925 & 0,900 & $-0{,}025$ & No \\
    Parcial~3      & 0,931 & 0,845 & $-0{,}086$ & Sí \\
    Prom.\ AMs     & 0,737 & 0,731 & $-0{,}006$ & No \\
    \bottomrule
  \end{tabular}
\end{table}

La Figura~\ref{fig:eras_comparacion} ilustra visualmente esta
comparación, tanto para las variables académicas como para los
patrones de engagement entre eras.

\begin{figure}[htbp]
  \centering
  \includegraphics[width=\textwidth]{09_comparacion_eras.png}
  \caption{\textbf{Izquierda:} AUC univariante de variables
           académicas por era pedagógica. La línea discontinua
           roja marca el umbral mínimo del SAT (AUC~$= 0{,}75$).
           \textbf{Derecha:} medianas de engagement por era
           y resultado académico, normalizadas para comparación
           entre variables de diferente escala.}
  \label{fig:eras_comparacion}
\end{figure}

% ............................................................
\subsubsection{Poder predictivo temprano del Parcial~1:
               tabla de riesgo empírico}
\label{sssec:riesgo_p1}
% ............................................................

El primer checkpoint del SAT (CP1, semana~6) opera únicamente
con la información disponible en ese momento, siendo el Parcial~1
el predictor individual más potente. La
Figura~\ref{fig:riesgo_p1} y la Tabla~\ref{tab:riesgo_p1}
cuantifican el riesgo empírico de reprobación por intervalo
de nota de P1 sobre \texttt{df\_activos} ($N = 554$).

\begin{figure}[htbp]
  \centering
  \includegraphics[width=\textwidth]{03_riesgo_p1.png}
  \caption{\textbf{Izquierda:} diagrama de dispersión entre la
           nota del Parcial~1 y la nota final, coloreado por
           resultado académico. Las líneas discontinuas marcan
           el umbral de aprobación (3{,}0). La correlación de
           Spearman se muestra en el recuadro.
           \textbf{Derecha:} tasa de reprobación empírica por
           intervalo de P1. La línea discontinua horizontal
           indica la tasa global de reprobación (25,7\,\%).}
  \label{fig:riesgo_p1}
\end{figure}

\begin{table}[htbp]
  \centering
  \caption{Tabla de riesgo empírico de reprobación por intervalo
           de nota en Parcial~1 (CP1, semana~6).}
  \label{tab:riesgo_p1}
  \begin{tabular}{ccrrrr}
    \toprule
    \textbf{Intervalo P1}
      & \textbf{Nivel}
      & \textbf{$N$}
      & \textbf{P(APROBÓ)}
      & \textbf{P(REPROBÓ)}
      & \textbf{$\bar{x}$ nota final} \\
    \midrule
    $[0{,}0;\;1{,}5)$  & Crítico  &  19 & 15,8\,\% & \textbf{84,2\,\%} & 2,16 \\
    $[1{,}5;\;2{,}5)$  & Alto     & 102 & 43,1\,\% & \textbf{56,9\,\%} & 2,80 \\
    $[2{,}5;\;3{,}0)$  & Moderado & 100 & 68,0\,\% &         32,0\,\%  & 3,10 \\
    $[3{,}0;\;3{,}5)$  & Límite   & 128 & 85,9\,\% &         14,1\,\%  & 3,44 \\
    $[3{,}5;\;4{,}0)$  & Seguro   &  95 & 83,2\,\% &         16,8\,\%  & 3,62 \\
    $[4{,}0;\;5{,}0]$  & Sólido   & 110 & 99,1\,\% &          0,9\,\%  & 4,18 \\
    \bottomrule
  \end{tabular}
\end{table}

Los resultados son contundentes en los extremos: estudiantes con
P1 $< 1{,}5$ tienen una probabilidad empírica del 84,2\,\% de
reprobar, mientras que con P1~$\geq 4{,}0$ esa probabilidad
cae al 0,9\,\%. Sin embargo, el intervalo $[2{,}5;\;3{,}5)$,
que concentra 228 estudiantes (41,2\,\% del total activo),
presenta tasas de reprobación intermedias (14--32\,\%) donde la
nota del parcial tiene menor poder de discriminación. Este es
precisamente el segmento donde el engagement y otros predictores
complementarios tienen mayor valor incremental.

% ............................................................
\subsubsection{Trayectorias académicas P1~$\to$~P2~$\to$~P3}
\label{sssec:trayectorias}
% ............................................................

Más allá del nivel absoluto de P1, la \emph{trayectoria}
de desempeño a lo largo de los tres parciales provee
información adicional sobre el resultado final. Se identifican
cinco patrones cualitativamente distintos
(Figura~\ref{fig:trayectorias} y Tabla~\ref{tab:trayectorias}).

\begin{figure}[htbp]
  \centering
  \includegraphics[width=\textwidth]{04_trayectorias_parciales.png}
  \caption{Trayectorias de desempeño académico entre parciales.
           \textbf{Izquierda:} frecuencia de cada trayectoria
           por resultado académico.
           \textbf{Derecha:} evolución de la nota promedio
           en cada parcial por trayectoria, con el porcentaje
           de aprobación de cada grupo en la leyenda.
           La línea discontinua marca el umbral de 3{,}0.}
  \label{fig:trayectorias}
\end{figure}

\begin{table}[htbp]
  \centering
  \caption{Patrones de trayectoria académica identificados en
           la secuencia P1~$\to$~P2~$\to$~P3. Los porcentajes
           se calculan dentro de cada grupo de trayectoria.}
  \label{tab:trayectorias}
  \begin{tabular}{lrrrr}
    \toprule
    \textbf{Trayectoria} & \textbf{Total}
      & \textbf{APROBÓ} & \textbf{REPROBÓ} & \textbf{\% APR.} \\
    \midrule
    Declive sostenido (P1~$>$~P2~$>$~P3)    &  40 &  22 &  18 & 55,0\,\% \\
    Estable   (P1~$\approx$~P2~$\approx$~P3)& 306 & 225 &  81 & 73,5\,\% \\
    Mejora sostenida (P1~$<$~P2~$<$~P3)     &  22 &  20 &   2 & 90,9\,\% \\
    Pico intermedio (P1~$<$~P2~$>$~P3)      & 129 &  99 &  30 & 76,7\,\% \\
    Recuperación post-P1 (P1~$<$~P2)        &  57 &  47 &  10 & 82,5\,\% \\
    \bottomrule
  \end{tabular}
\end{table}

El grupo de \textbf{declive sostenido} concentra la mayor tasa
de reprobación (45,0\,\%) y es identificable a partir de la
comparación P1~$\to$~P2 disponible en CP2. La
\textbf{mejora sostenida} es el patrón con mayor probabilidad
de aprobación (90,9\,\%), lo que indica que la trayectoria
ascendente es una señal de éxito más fiable que el nivel
absoluto de P1.

Más relevante aún para el diseño del SAT: 115 de los 221
estudiantes con P1~$< 3{,}0$ terminaron aprobando (52,0\,\%),
con nota media de Parcial~2 $= 3{,}46$ y nota final media
$= 3{,}38$. Este fenómeno de \textit{recuperación académica}
demuestra que un clasificador basado únicamente en P1 tendría
una tasa de falsas alarmas del 52\,\% entre los estudiantes de
riesgo moderado, motivando explícitamente la revisión del estado
de alerta en el segundo checkpoint (CP2, semana~11).

% ------------------------------------------------------------
\subsection{Análisis de la Actividad en Plataforma
            (Engagement)}
\label{subsec:eng}
% ------------------------------------------------------------

El análisis de engagement se realiza sobre
\texttt{df\_activos\_engagement} ($N = 502$): 373 aprobados y
129 reprobados. La Figura~\ref{fig:composicion_modelado}
muestra la composición de este subconjunto.

\begin{figure}[htbp]
  \centering
  \includegraphics[width=0.92\textwidth]{12_composicion_modelado.png}
  \caption{Composición del subconjunto de modelado
           (\texttt{df\_activos\_engagement}, $N = 502$).
           \textbf{Izquierda:} distribución de clases. El
           desbalance 74,3\,\%~/ 25,7\,\% motiva el uso de
           técnicas de balanceo (SMOTE, SMOTEENN).
           \textbf{Derecha:} descomposición por semestre, con
           el porcentaje de reprobados sobre cada barra.}
  \label{fig:composicion_modelado}
\end{figure}

% ............................................................
\subsubsection{Descripción de las métricas de engagement}
\label{sssec:vars_eng}
% ............................................................

Bloque Neón registra el comportamiento de estudio en dos
dimensiones complementarias:

\begin{itemize}
  \item \textbf{Amplitud de estudio (cobertura):}
    módulos únicos visitados (\texttt{parcial\_X\_modulos\_unicos}),
    temas únicos consultados (\texttt{parcial\_X\_temas\_unicos})
    y visitas totales al contenido (\texttt{parcial\_X\_visitas}).

  \item \textbf{Profundidad de estudio (intensidad):}
    tiempo promedio por visita
    (\texttt{parcial\_X\_tiempo\_por\_visita}) y visitas promedio
    por tema (\texttt{parcial\_X\_visitas\_por\_tema}).
\end{itemize}

Además, se dispone del tiempo acumulado hasta cada checkpoint
(\texttt{engagement\_hasta\_p1},
\texttt{engagement\_hasta\_p2}) y del total del semestre
(\texttt{total\_tiempo\_hrs}).

% ............................................................
\subsubsection{Diferencias estadísticas entre aprobados y
               reprobados: prueba de Mann-Whitney}
\label{sssec:mann_eng}
% ............................................................

La Figura~\ref{fig:engagement_boxplots} presenta las
distribuciones de las principales variables de engagement
disponibles en CP1 por semestre y resultado, y la
Tabla~\ref{tab:mw_engagement} reporta los estadísticos
completos de las pruebas Mann-Whitney.

\begin{figure}[htbp]
  \centering
  \includegraphics[width=\textwidth]{07_engagement_boxplots.png}
  \caption{Distribución de las cuatro principales variables de
           engagement en CP1 por semestre y resultado académico.
           El sombreado naranja indica la era antigua y el verde
           la nueva. En cada subgráfico se reportan el $p$-valor
           de Mann-Whitney, el AUC univariante y el coeficiente
           de rango biserial $r_b$.}
  \label{fig:engagement_boxplots}
\end{figure}

\begin{table}[htbp]
  \centering
  \caption{Prueba Mann-Whitney para variables de engagement
           disponibles en CP1. $\tilde{x}$: mediana; $r_b$:
           coeficiente de rango biserial \cite{Kerby2014}.
           Significancia: *** $p < 0{,}001$; * $p < 0{,}05$.}
  \label{tab:mw_engagement}
  \begin{tabular}{lrrrrc}
    \toprule
    \textbf{Variable}
      & \textbf{$\tilde{x}$ APR.}
      & \textbf{$\tilde{x}$ REP.}
      & \textbf{AUC}
      & \textbf{$r_b$}
      & \textbf{$p$-valor} \\
    \midrule
    Tiempo hasta P1 (h)
      & 8,27  & 5,71  & 0,602 & 0,204
      & $2{,}8\times10^{-4}$ *** \\
    Módulos únicos P1
      & 7,00  & 6,00  & 0,601 & 0,203
      & $2{,}7\times10^{-4}$ *** \\
    Visitas contenido P1
      & 37,00 & 32,00 & 0,575 & 0,150
      & $1{,}1\times10^{-2}$ *   \\
    Temas únicos P1
      & 12,00 & 11,00 & 0,567 & 0,134
      & $2{,}3\times10^{-2}$ *   \\
    Visitas por tema P1
      & 3,20  & 2,83  & 0,565 & 0,130
      & $2{,}8\times10^{-2}$ *   \\
    Minutos por visita P1
      & 12,26 & 11,07 & 0,562 & 0,125
      & $3{,}5\times10^{-2}$ *   \\
    \bottomrule
  \end{tabular}
\end{table}

Todas las variables de engagement en CP1 presentan diferencias
estadísticamente significativas ($p < 0{,}05$), pero los
tamaños del efecto son \textbf{pequeños}
($r_b \in [0{,}12;\;0{,}20]$), muy por debajo de los observados
para las variables académicas (Tabla~\ref{tab:mw_academicas}).

\begin{table}[htbp]
  \centering
  \caption{Prueba Mann-Whitney para variables académicas.
           Todas presentan tamaños del efecto grandes
           ($d_{\text{Cohen}} \geq 0{,}91$)
           con significancia $p < 10^{-15}$.}
  \label{tab:mw_academicas}
  \begin{tabular}{lrrrrrr}
    \toprule
    \textbf{Variable}
      & \textbf{$\tilde{x}$ APR.}
      & \textbf{$\tilde{x}$ REP.}
      & \textbf{AUC}
      & \textbf{$r_b$}
      & \textbf{$d_{\text{Cohen}}$}
      & \textbf{$p$-valor} \\
    \midrule
    Parcial~1       & 3,42 & 2,43 & 0,810 & 0,620 & 1,24
      & $1{,}8\times10^{-28}$ \\
    Parcial~2       & 3,80 & 2,33 & 0,861 & 0,722 & 1,56
      & $6{,}7\times10^{-38}$ \\
    Parcial~3       & 3,48 & 1,70 & 0,882 & 0,763 & 1,66
      & $4{,}7\times10^{-42}$ \\
    Prom.\ AMs      & 4,11 & 3,12 & 0,728 & 0,456 & 0,91
      & $2{,}9\times10^{-16}$ \\
    Prom.\ Quizzes  & 4,00 & 2,91 & 0,821 & 0,642 & 1,36
      & $1{,}0\times10^{-20}$ \\
    \bottomrule
  \end{tabular}
\end{table}

% ............................................................
\subsubsection{Poder discriminativo individual: ranking de AUC}
\label{sssec:auc_ranking}
% ............................................................

La Figura~\ref{fig:auc_univariante} resume el poder discriminativo
individual de cada variable disponible en CP1 mediante su AUC
univariante, ordenado de mayor a menor.

\begin{figure}[htbp]
  \centering
  \includegraphics[width=0.92\textwidth]{06_auc_univariante_cp1.png}
  \caption{AUC univariante de cada predictor candidato para CP1.
           Las variables académicas (azul) dominan ampliamente;
           las de engagement (verde) se ubican en la zona
           $[0{,}56;\;0{,}60]$, superando la clasificación
           aleatoria pero por debajo del umbral mínimo del SAT
           (0{,}75, línea verde discontinua). Prom.\ Quizzes no
           está disponible en CP1 para todos los semestres.}
  \label{fig:auc_univariante}
\end{figure}

La brecha entre el AUC de Parcial~1 (0,810) y el del mejor
predictor de engagement (Tiempo hasta P1: 0,602) es de 20,8
puntos porcentuales. Esto confirma que el engagement es un
predictor \emph{complementario}, no un sustituto de las
calificaciones. Su valor radica en los casos marginales donde
la nota de P1 tiene menor poder de discriminación
(zona $[2{,}5;\;3{,}5)$ de la tabla de riesgo).

% ............................................................
\subsubsection{Valor incremental del engagement y correlación
               parcial}
\label{sssec:eng_incremental}
% ............................................................

Para cuantificar el aporte del engagement \emph{más allá} de lo
que ya captura Parcial~1, se calcula la correlación de Spearman
entre el tiempo de engagement hasta P1 y la nota final,
controlando por la nota de P1. El resultado es
$\rho_{\text{parcial}} = 0{,}166$ ($p = 1{,}9 \times 10^{-4}$,
$N = 502$), indicando que el engagement explica un 2--3\,\% de
varianza adicional en el resultado. Aunque modesto en términos
absolutos, este aporte es estadísticamente robusto y
operacionalmente relevante en la zona de incertidumbre del
clasificador.

% ............................................................
\subsubsection{Patrones de acumulación de engagement entre
               checkpoints}
\label{sssec:eng_evolucion}
% ............................................................

La Figura~\ref{fig:engagement_evolucion} explora cómo se acumula
el engagement entre los dos checkpoints del SAT.

\begin{figure}[htbp]
  \centering
  \includegraphics[width=\textwidth]{10_evolucion_engagement.png}
  \caption{\textbf{Izquierda:} correlación entre el tiempo de
           engagement hasta P1 y hasta P2 por resultado académico
           ($\rho_{\text{Spearman}}$ reportado). La diagonal
           discontinua indica el punto de no-incremento entre
           checkpoints.
           \textbf{Derecha:} distribución del incremento de
           engagement en cada período por resultado (violines).
           Los $p$-valores corresponden a pruebas Mann-Whitney
           entre estados.}
  \label{fig:engagement_evolucion}
\end{figure}

La correlación entre el engagement hasta P1 y hasta P2 es alta
para ambos grupos (aprobados: $\rho \approx 0{,}86$;
reprobados: similar), lo que indica \textbf{persistencia de los
hábitos de estudio}: los estudiantes que estudian más en el primer
período tienden a hacerlo también en el segundo. Los aprobados
acumulan en el período P1~$\to$~P2 una mediana de 6,80 horas
adicionales frente a 4,85 de los reprobados; esta diferencia
se amplía hacia el final del semestre (5,29~h vs.~3,28~h en
el período P2~$\to$~fin). Este patrón sugiere que los estudiantes
en riesgo no solo estudian menos, sino que además reducen
progresivamente su dedicación conforme avanza el semestre.

% ------------------------------------------------------------
\subsection{Análisis de Correlaciones y Multicolinealidad}
\label{subsec:corr}
% ------------------------------------------------------------

% ............................................................
\subsubsection{Correlaciones entre variables de engagement}
\label{sssec:corr_eng}
% ............................................................

La Figura~\ref{fig:corr_engagement} presenta la matriz de
correlaciones de Spearman entre las variables de engagement y
la dispersión entre tiempo hasta P1 y nota del Parcial~1.

\begin{figure}[htbp]
  \centering
  \includegraphics[width=\textwidth]{08_correlaciones_engagement.png}
  \caption{\textbf{Izquierda:} matriz de correlaciones de Spearman
           entre variables de engagement (triángulo inferior).
           Las celdas en rojo oscuro indican correlaciones
           altas ($r > 0{,}80$).
           \textbf{Derecha:} dispersión entre tiempo de
           engagement hasta P1 y nota del Parcial~1 por resultado
           académico, con línea de regresión y correlación de
           Spearman.}
  \label{fig:corr_engagement}
\end{figure}

Las correlaciones intra-período son altas: los pares más
problemáticos son tiempo hasta P1 $\leftrightarrow$ tiempo
hasta P2 ($r = 0{,}864$), temas~P1 $\leftrightarrow$
módulos~P1 ($r = 0{,}861$), tiempo~P2 $\leftrightarrow$
tiempo total ($r = 0{,}852$) y temas~total $\leftrightarrow$
módulos~total ($r = 0{,}941$). Esta multicolinealidad no
impide el modelado pero requiere regularización en los
modelos lineales y justifica la selección de subconjuntos
representativos de cada dimensión de engagement en lugar
de incluir todas las variables simultáneamente.

\subsubsection{Factor de inflación de varianza (VIF)}
\label{sssec:vif}

El análisis VIF sobre el conjunto de predictores candidatos
revela que las variables académicas presentan los valores más
altos (Parcial~1: VIF~$= 13{,}7$; AM~2: VIF~$= 10{,}4$;
AM~1: VIF~$= 9{,}4$), mientras que las de engagement son más
moderadas (módulos únicos P1: VIF~$= 7{,}2$;
tiempo hasta P1: VIF~$= 4{,}0$;
tiempo por visita P1: VIF~$= 3{,}9$). Ningún valor supera el
umbral de 15 que produce inestabilidad severa, pero la magnitud
de algunos VIF justifica el uso de regularización L1/L2 en los
modelos de regresión logística (ver Sección~\ref{sec:metodologia}).

% ------------------------------------------------------------
\subsection{Síntesis: Decisiones de Diseño del Dataset}
\label{subsec:sintesis_datos}
% ------------------------------------------------------------

La Tabla~\ref{tab:resumen_decisiones} sintetiza las decisiones
de preprocesamiento tomadas a partir del análisis exploratorio.

\begin{table}[htbp]
  \centering
  \caption{Resumen de decisiones de diseño del dataset con
           justificación empírica e impacto cuantificado.}
  \label{tab:resumen_decisiones}
  \small
  \begin{tabular}{p{3.5cm}p{5.0cm}p{4.8cm}}
    \toprule
    \textbf{Decisión} & \textbf{Justificación} & \textbf{Impacto} \\
    \midrule
    Excluir RETIRADO y SIN~NOTA
      & Etiqueta no observable; sesgo de retiro endógeno;
        heterogeneidad semántica
      & $N$: 764~$\to$~554 \\[3pt]
    Excluir estudiantes sin engagement
      & Cero ambiguo (75\,\% en semestre de arranque);
        $\Delta F_1 < 0{,}02$ al excluirlos
      & $N$: 554~$\to$~502 \\[3pt]
    Excluir variables de actividades evaluativas
      & AUC~$< 0{,}50$; $p > 0{,}05$ (sin discriminación)
      & $-3$ predictores candidatos \\[3pt]
    Incluir \texttt{era\_encoded}
      & $\Delta$AUC(P1) $= -0{,}087$;
        engagement mediano cae 60\,\% en era nueva
      & Control en todos los modelos \\[3pt]
    Incluir \texttt{modalidad\_am}
      & AM-3 grupal vs.\ individual:
        $p = 1{,}4\times10^{-9}$, $d = 0{,}40$
      & Control en modelos con AMs \\[3pt]
    Incluir engagement a pesar de efecto pequeño
      & $\rho_{\text{parcial}} = 0{,}166$, $p < 0{,}001$;
        valor incremental en zona marginal de P1
      & $+6$ predictores en CP1 \\[3pt]
    Aplicar balanceo SMOTE / SMOTEENN
      & 74,3\,\% APROBÓ vs.\ 25,7\,\% REPROBÓ;
        Recall sin balanceo $\approx 0{,}55$
      & Solo en training fold; nunca en evaluación \\
    \bottomrule
  \end{tabular}
\end{table}

El dataset de modelado resultante (\texttt{df\_activos\_engagement},
$N = 502$) presenta tres características estructurales que
condicionan el diseño metodológico de la siguiente sección:
(i)~\textbf{desbalance de clases} (74,3\,\% vs.\ 25,7\,\%),
que exige técnicas de remuestreo para evitar que los modelos
sesguen la predicción hacia la clase mayoritaria;
(ii)~\textbf{estructura longitudinal multisemestral}, que
invalida la validación cruzada aleatoria estándar y exige el
esquema LOSO; y (iii)~\textbf{heterogeneidad de era pedagógica},
que requiere variables de control y análisis estratificado.
Estos tres aspectos se abordan en detalle en la sección de
metodología.

% ============================================================
%  REFERENCIAS
% ============================================================

\begin{thebibliography}{9}

\bibitem{Lopez2021}
J.~López-Zambrano, J.~A. Lara~Torralbo, y C.~Romero,
``Early prediction of student learning performance through data
mining: A systematic review,''
\textit{Psicothema}, vol.~33, no.~3, pp.~456--465, 2021.

\bibitem{Tinto1993}
V.~Tinto,
\textit{Leaving College: Rethinking the Causes and Cures of
Student Attrition}, 2.ª~ed.
Chicago, IL, USA: University of Chicago Press, 1993.

\bibitem{Kerby2014}
D.~S. Kerby,
``The simple difference formula: An approach to teaching
nonparametric correlation,''
\textit{Compr. Psychol.}, vol.~3, p.~11.IT.3.4.1, 2014.

\bibitem{He2009}
H.~He y E.~A. Garcia,
``Learning from imbalanced data,''
\textit{IEEE Trans. Knowl. Data Eng.}, vol.~21, no.~9,
pp.~1263--1284, sep. 2009.

\end{thebibliography}

\end{document}
 y eliminar el
%  preámbulo y los comandos \begin{document}/\end{document}.
%  Cambiar [english] a [spanish] con texlive-lang-spanish.
%  Las figuras se cargan desde la carpeta figures/ relativa.
% ============================================================

\documentclass[12pt, letterpaper]{article}

\usepackage[utf8]{inputenc}
\usepackage[T1]{fontenc}
\usepackage[english]{babel}
\usepackage{amsmath, amssymb}
\usepackage{geometry}
\usepackage{setspace}
\usepackage{booktabs}
\usepackage{multirow}
\usepackage{array}
\usepackage{graphicx}
\usepackage{subcaption}
\usepackage{xcolor}
\usepackage{cite}
\usepackage[hidelinks]{hyperref}
\usepackage{url}
\usepackage{rotating}

\geometry{top=2.5cm, bottom=2.5cm, left=3cm, right=2.5cm}
\onehalfspacing
\renewcommand{\refname}{Referencias}

\graphicspath{{figures/}}

% ============================================================
\begin{document}
% ============================================================

\section{Descripción y Exploración de Datos}
\label{sec:exploracion}

Este capítulo caracteriza el conjunto de datos empleado en el SAT,
documenta las decisiones de preprocesamiento adoptadas y presenta
los hallazgos del análisis exploratorio que fundamentan el diseño
del sistema de modelado. La presentación se organiza en seis
bloques: (i)~fuente y composición del dataset;
(ii)~definición y justificación de los subconjuntos de trabajo;
(iii)~calidad y estructura de los datos;
(iv)~análisis estadístico de las notas académicas;
(v)~análisis de la actividad en plataforma; y
(vi)~análisis de correlaciones y síntesis de decisiones de diseño.

% ------------------------------------------------------------
\subsection{Fuente y Composición del Conjunto de Datos}
\label{subsec:fuente}
% ------------------------------------------------------------

Los datos provienen del curso IIND-2104 \textit{Modelos
Probabilísticos} de la Universidad de los Andes, dictado entre
los semestres 202320 y 202520 —cinco cohortes consecutivas
cubiertas entre 2023 y 2025. La información se integra a partir
de dos fuentes complementarias:

\begin{itemize}
  \item \textbf{Sistema de gestión académica institucional:}
    calificaciones oficiales de los tres parciales y de las
    evaluaciones continuas (actividades modulares, quizzes,
    talleres y proyectos, según el formato de cada semestre).

  \item \textbf{Plataforma Bloque Neón:} sistema de gestión del
    aprendizaje (LMS) que registra, por sesión y por módulo
    de contenido, el tiempo de estudio, los temas consultados,
    los módulos visitados y la frecuencia de acceso de cada
    estudiante.
\end{itemize}

El dataset crudo consolidado cuenta con \textbf{764~registros y
77~variables}, con cuatro posibles resultados académicos finales:
\texttt{APROBÓ} (413 estudiantes, 54,1\,\%),
\texttt{REPROBÓ} (141, 18,5\,\%),
\texttt{RETIRADO} (164, 21,5\,\%) y
\texttt{SIN~NOTA} (46, 6,0\,\%).

La Figura~\ref{fig:composicion} ilustra cómo se distribuyen
estos resultados por semestre en términos absolutos y
porcentuales, y la Tabla~\ref{tab:composicion_semestral} detalla
las frecuencias exactas.

\begin{figure}[htbp]
  \centering
  \includegraphics[width=\textwidth]{01_composicion_dataset.png}
  \caption{Composición del dataset por semestre y resultado
           académico. \textbf{Izquierda:} distribución absoluta.
           \textbf{Derecha:} distribución porcentual con
           etiquetas de proporción. El sombreado naranja
           corresponde a la era pedagógica antigua (202320--202410)
           y el verde a la nueva (202420--202520).}
  \label{fig:composicion}
\end{figure}

\begin{table}[htbp]
  \centering
  \caption{Composición por semestre, era pedagógica y resultado
           académico final.}
  \label{tab:composicion_semestral}
  \small
  \begin{tabular}{lccrrrrrrr}
    \toprule
    \textbf{Semestre} & \textbf{Era} & \textbf{Total}
      & \multicolumn{2}{c}{\textbf{APROBÓ}}
      & \multicolumn{2}{c}{\textbf{REPROBÓ}}
      & \multicolumn{2}{c}{\textbf{RETIRADO}}
      & \textbf{SIN} \\
    \cmidrule(lr){4-5}\cmidrule(lr){6-7}\cmidrule(lr){8-9}
    & & & $n$ & \% & $n$ & \% & $n$ & \% & \textbf{NOTA} \\
    \midrule
    202320 & Antigua & 170 & 116 & 68,2 &  27 & 15,9 &  27 & 15,9 & 0  \\
    202410 & Antigua & 150 &  68 & 45,3 &  14 &  9,3 &  30 & 20,0 & 38 \\
    202420 & Nueva   & 164 &  90 & 54,9 &  33 & 20,1 &  40 & 24,4 & 1  \\
    202510 & Nueva   & 140 &  70 & 50,0 &  47 & 33,6 &  21 & 15,0 & 2  \\
    202520 & Nueva   & 140 &  69 & 49,3 &  20 & 14,3 &  46 & 32,9 & 5  \\
    \midrule
    \textbf{Total} & — & \textbf{764}
      & \textbf{413} & \textbf{54,1}
      & \textbf{141} & \textbf{18,5}
      & \textbf{164} & \textbf{21,5}
      & \textbf{46} \\
    \bottomrule
  \end{tabular}
\end{table}

Dos fenómenos merecen atención especial. Primero, la tasa de
reprobación en 202510 (33,6\,\%) es la más alta del período y
más del doble de la de 202410 (9,3\,\%), lo que sugiere
variabilidad intercohorte que el modelo debe generalizar. Segundo,
el semestre 202520 registra la mayor tasa de retiro (32,9\,\%),
fenómeno atípico que podría obedecer a cambios en la política
institucional de permanencia o a condiciones externas no
capturadas en los datos.

% ............................................................
\subsubsection{Dos eras pedagógicas}
\label{sssec:eras}
% ............................................................

Un rasgo estructural del dataset es la coexistencia de dos
\textbf{formatos de evaluación continua}:

\begin{itemize}
  \item \textbf{Era Antigua} (202320 y 202410; $n = 320$ estudiantes
    activos): las evaluaciones continuas consisten en talleres
    individuales y proyectos por fases grupales.

  \item \textbf{Era Nueva} (202420, 202510 y 202520; $n = 444$
    estudiantes activos): las evaluaciones continuas son
    actividades modulares (AMs) y quizzes individuales.
\end{itemize}

Dentro de la era nueva existe una distinción adicional: en
202420 las AMs se realizaron \textbf{en grupo}, mientras que
a partir de 202510 pasaron a ser \textbf{individuales}. Este
cambio produce diferencias estadísticamente significativas en la
distribución de AM-3 ($U$-test, $p = 1{,}39 \times 10^{-9}$,
$d_{\text{Cohen}} = 0{,}40$, AUC$_{\text{grupo}} = 0{,}369$
vs.\ AUC$_{\text{ind.}} = 0{,}246$), lo que justifica incluir la
variable \texttt{modalidad\_am} como control en los modelos
entrenados con datos del formato nuevo.

% ------------------------------------------------------------
\subsection{Definición y Justificación de los Subconjuntos
            de Trabajo}
\label{subsec:subsets}
% ------------------------------------------------------------

El dataset bruto requiere dos depuraciones sucesivas antes de
ser apto para el modelado de clasificación.

% ............................................................
\subsubsection{Exclusión de estudiantes retirados y sin
               calificación}
\label{sssec:excl_retirados}
% ............................................................

Los 164 estudiantes con resultado \texttt{RETIRADO} y los
46 con \texttt{SIN~NOTA} se excluyen por tres razones:

\begin{enumerate}
  \item \textbf{Ausencia de etiqueta objetivo observable.}
    El SAT modela $P(\text{REPROBÓ} \mid \mathbf{x}_i)$, la
    probabilidad de que un estudiante \emph{finalice} el semestre
    reprobando. Los estudiantes retirados abandonaron el curso
    antes de producirse el resultado; su etiqueta final es
    desconocida y no puede inferirse sin supuestos fuertes.

  \item \textbf{Sesgo de retiro endógeno.}
    La decisión de retirarse está correlacionada con el
    desempeño anticipado: los estudiantes que perciben que van
    a reprobar se retiran en mayor proporción \cite{Tinto1993}.
    Incluirlos en la clase \texttt{REPROBÓ} introduciría sesgo
    de selección en las variables de engagement: un estudiante
    retirado en la semana~4 tiene engagement mínimo no por falta
    de habilidad sino por abandono temprano, contaminando el
    predictor más utilizado.

  \item \textbf{Heterogeneidad semántica.}
    La predicción del retiro como fenómeno independiente
    requeriría modelos multi-clase y supuestos adicionales
    sobre su equivalencia con la reprobación, lo cual está
    fuera del alcance del presente trabajo.
\end{enumerate}

El subconjunto resultante es \textbf{\texttt{df\_activos}}
($N = 554$; APROBÓ: 413, REPROBÓ: 141).

% ............................................................
\subsubsection{Exclusión adicional de estudiantes sin actividad
               en plataforma}
\label{sssec:excl_engagement}
% ............................................................

Del subconjunto \texttt{df\_activos}, 52 estudiantes (9,4\,\%)
presentan \textbf{engagement total nulo} en Bloque Neón durante
todo el semestre (i.e., \texttt{engagement\_hasta\_p1} $= 0$).
Estos casos se excluyen del subconjunto de modelado por las
razones siguientes:

\begin{enumerate}
  \item \textbf{Ambigüedad del cero.}
    Existen dos interpretaciones mutuamente excluyentes:
    (a)~el estudiante genuinamente no utilizó la plataforma,
    o (b)~ocurrió un error en el pipeline de integración de
    datos. Ambas hipótesis son plausibles: 39 de los 52 casos
    (75\,\%) se concentran en el semestre 202320, el primero
    en incorporar Bloque Neón, cuando el sistema estaba aún
    en fase de implementación. No es posible distinguir entre
    causas con la información disponible.

  \item \textbf{Riesgo de imputación espuria.}
    Imputar cero como valor real de engagement asumiría que
    ``no usar la plataforma'' es una señal homogénea entre
    semestres, cuando los patrones de no-uso varían
    considerablemente entre cohortes.

  \item \textbf{Impacto estadístico marginal confirmado.}
    Se realizó un análisis de sensibilidad: los modelos
    entrenados sobre \texttt{df\_activos} completo ($N = 554$)
    y sobre el subconjunto filtrado ($N = 502$) arrojan
    diferencias en F1 menores a $\Delta F_1 < 0{,}02$,
    confirmando que la exclusión no altera las conclusiones.
\end{enumerate}

El subconjunto de modelado final es
\textbf{\texttt{df\_activos\_engagement}}
($N = 502$; APROBÓ:~373, REPROBÓ:~129).
La Tabla~\ref{tab:subsets} resume los tres subconjuntos
empleados en el análisis.

\begin{table}[htbp]
  \centering
  \caption{Subconjuntos de trabajo. \texttt{df\_activos\_engagement}
           es el conjunto empleado para el entrenamiento y
           evaluación de todos los modelos del SAT.}
  \label{tab:subsets}
  \begin{tabular}{lrrrl}
    \toprule
    \textbf{Subconjunto} & \textbf{$N$}
      & \textbf{APROBÓ} & \textbf{REPROBÓ}
      & \textbf{Uso principal} \\
    \midrule
    \texttt{df\_all}
      & 764 & 413 & 141
      & Tasas globales, análisis de retiro \\
    \texttt{df\_activos}
      & 554 & 413 & 141
      & Análisis de notas y distribuciones \\
    \texttt{df\_activos\_engagement}
      & \textbf{502} & \textbf{373} & \textbf{129}
      & \textbf{Entrenamiento y evaluación de modelos} \\
    \bottomrule
  \end{tabular}
\end{table}

% ------------------------------------------------------------
\subsection{Calidad y Estructura de los Datos}
\label{subsec:calidad}
% ------------------------------------------------------------

% ............................................................
\subsubsection{Nulos estructurales versus problemáticos}
\label{sssec:nulos}
% ............................................................

El dataset presenta un volumen considerable de valores faltantes,
en su mayoría atribuibles al cambio de formato pedagógico.
Se distinguen dos categorías.

\paragraph{Nulos estructurales (esperados).}
Son variables que solo existen en una era y, por tanto, están
ausentes de forma legítima en la otra:
\texttt{quiz\_4/5/6} (77,8\,\% de nulos, solo en algunos
semestres del formato nuevo),
\texttt{taller\_1/2/3} (59,4\,\%, exclusivo de la era antigua),
\texttt{fase\_1/2} (59,4\,\%, proyectos grupales de la era antigua),
y \texttt{am\_4/5/6} (56,7\,\%, AMs opcionales no presentes en
todos los semestres). Estos nulos no representan problemas de
calidad sino de diseño del instrumento de evaluación.

\paragraph{Nulos problemáticos (casos especiales).}
Son faltantes en variables que deberían estar presentes para
todos los estudiantes activos:
\begin{itemize}
  \item \textbf{Parcial~1 $= 0$ ($n = 1$):} un único estudiante
    registra nota cero. No es posible distinguir ausencia de nota
    real. Se retiene con variable indicadora
    \texttt{no\_presento\_p1}.
  \item \textbf{Parcial~3 $= 0$ ($n = 15$; 13 reprobaron):}
    concentrado en estudiantes encaminados a reprobar, consistente
    con ausentismo por desmotivación al final del semestre.
    Se retienen con su valor real.
  \item \textbf{Notas superiores a 5{,}0:} presentes por
    bonificaciones otorgadas por docentes. Son valores legítimos
    y se retienen sin transformación.
\end{itemize}

% ............................................................
\subsubsection{Variables de plataforma sin poder discriminativo}
\label{sssec:excl_vars}
% ............................................................

La plataforma Bloque Neón registra tanto métricas de
\textit{estudio de contenidos} como métricas de
\textit{actividades evaluativas en línea}. El análisis estadístico
univariante sobre estas últimas —\texttt{actividades\_n\_realizadas},
\texttt{actividades\_tiempo\_total} y
\texttt{actividades\_visitas\_total}— revela ausencia de
poder discriminativo entre aprobados y reprobados
(Tabla~\ref{tab:excl_vars}), lo que se ilustra también en la
Figura~\ref{fig:vars_excluidas}.

\begin{table}[htbp]
  \centering
  \caption{Variables de actividades evaluativas de la plataforma.
           AUC~$\approx 0{,}5$ y $p > 0{,}05$ indican que no
           discriminan mejor que la clasificación aleatoria.}
  \label{tab:excl_vars}
  \begin{tabular}{lrrrr}
    \toprule
    \textbf{Variable}
      & \textbf{Med.\ APROBÓ}
      & \textbf{Med.\ REPROBÓ}
      & \textbf{AUC}
      & \textbf{$p$-valor (MW)} \\
    \midrule
    \texttt{actividades\_n\_realizadas}   & 1,00 & 1,00 & 0,473 & 0,830 \\
    \texttt{actividades\_tiempo\_total}   & 1,15 & 1,10 & 0,498 & 0,532 \\
    \texttt{actividades\_visitas\_total}  & 9,00 & 11,00 & 0,465 & 0,886 \\
    \bottomrule
  \end{tabular}
\end{table}

\begin{figure}[htbp]
  \centering
  \includegraphics[width=\textwidth]{11_variables_excluidas.png}
  \caption{Distribuciones de las tres variables de actividades
           evaluativas para aprobados (azul) y reprobados (rojo).
           Los valores de $p > 0{,}05$ y AUC~$< 0{,}5$
           confirman la ausencia de poder discriminativo,
           justificando su exclusión del conjunto de predictores.}
  \label{fig:vars_excluidas}
\end{figure}

Nótese que, para \texttt{actividades\_visitas\_total}, la mediana
de los reprobados (11) supera la de los aprobados (9), con
AUC~$= 0{,}465$ —peor que la clasificación aleatoria—, lo que
sugiere que los estudiantes con mayor número de visitas a
actividades no necesariamente obtienen mejores resultados. Las
tres variables son excluidas del conjunto de predictores para
evitar introducir ruido en los modelos.

% ------------------------------------------------------------
\subsection{Análisis Estadístico de las Notas Académicas}
\label{subsec:notas}
% ------------------------------------------------------------

El análisis de notas se realiza sobre \texttt{df\_activos}
($N = 554$) para preservar la distribución real sin los
filtros de engagement.

% ............................................................
\subsubsection{Distribución de la nota final}
\label{sssec:dist_nota}
% ............................................................

La nota final (escala 0--5) presenta una distribución bimodal
con dos concentraciones claramente separadas: una en torno a
2,6 (zona de reprobación) y otra en torno a 3,7 (zona de
aprobación cómoda), como se observa en la
Figura~\ref{fig:dist_nota}.

\begin{figure}[htbp]
  \centering
  \includegraphics[width=\textwidth]{02_distribucion_nota_final.png}
  \caption{Distribución de la nota final.
           \textbf{Izquierda:} histogramas superpuestos con
           curva de densidad KDE por resultado académico.
           La línea discontinua vertical indica el umbral de
           aprobación (3{,}0).
           \textbf{Derecha:} diagramas de caja por semestre y
           resultado; el sombreado distingue las eras pedagógicas.}
  \label{fig:dist_nota}
\end{figure}

Las estadísticas descriptivas detalladas se presentan en la
Tabla~\ref{tab:desc_nota}. La prueba de Mann-Whitney confirma
la separación altamente significativa entre grupos
($U = 58\,233{,}0$, $p = 9{,}97 \times 10^{-71}$,
$d_{\text{Cohen}} = 2{,}41$). Esta magnitud del efecto implica
que la distribución de notas de los aprobados y la de los
reprobados se solapan apenas en una pequeña región central
(alrededor de $[2{,}8;\;3{,}2]$), que es precisamente la zona
donde los modelos de clasificación concentran su incertidumbre.

\begin{table}[htbp]
  \centering
  \caption{Estadísticas descriptivas de la nota final por resultado
           y era pedagógica. $\bar{x}$: media; $\tilde{x}$: mediana;
           $s$: desviación estándar.}
  \label{tab:desc_nota}
  \begin{tabular}{lrrrr}
    \toprule
    \textbf{Grupo} & \textbf{$N$}
      & \textbf{$\bar{x}$} & \textbf{$\tilde{x}$} & \textbf{$s$} \\
    \midrule
    APROBÓ (global)                & 413 & 3,724 & 3,620 & 0,527 \\
    REPROBÓ (global)               & 141 & 2,418 & 2,632 & 0,583 \\
    \midrule
    Era Antigua — APROBÓ           & 184 & 3,777 & 3,680 & 0,504 \\
    Era Antigua — REPROBÓ          &  41 & 2,446 & 2,720 & 0,671 \\
    Era Nueva   — APROBÓ           & 229 & 3,682 & 3,530 & 0,541 \\
    Era Nueva   — REPROBÓ          & 100 & 2,407 & 2,605 & 0,545 \\
    \bottomrule
  \end{tabular}
\end{table}

% ............................................................
\subsubsection{Poder predictivo de los parciales:
               análisis por era pedagógica}
\label{sssec:parciales_era}
% ............................................................

La Figura~\ref{fig:parciales_era} compara la distribución de
las tres calificaciones de parciales entre aprobados y reprobados,
estratificada por era pedagógica.

\begin{figure}[htbp]
  \centering
  \includegraphics[width=\textwidth]{05_parciales_era_resultado.png}
  \caption{Distribuciones de Parcial~1, Parcial~2 y Parcial~3
           por resultado académico, estratificadas por era
           pedagógica. Cada subgráfico incluye el $p$-valor
           de la prueba Mann-Whitney, el AUC univariante y el
           coeficiente de rango biserial~$r_b$.}
  \label{fig:parciales_era}
\end{figure}

La Tabla~\ref{tab:auc_era} sintetiza los AUC univariantes por
variable y era. El hallazgo más relevante es la \textbf{reducción
del poder predictivo de P1 en la era nueva}
($\Delta\text{AUC} = -0{,}087$). Esta caída puede atribuirse
al mayor peso relativo de las actividades continuas (quizzes y
AMs) en el nuevo currículo, que diluye la señal del primer
examen. Contrariamente, el Parcial~2 mantiene una capacidad
predictiva alta y estable en ambas eras
($\Delta\text{AUC} = -0{,}025$), lo que refuerza el valor
del segundo checkpoint (CP2) del SAT.

\begin{table}[htbp]
  \centering
  \caption{AUC univariante de las variables académicas por era
           pedagógica. $\Delta$AUC indica la variación entre eras.
           Diferencias $> 5\,\%$ se consideran sustanciales.}
  \label{tab:auc_era}
  \begin{tabular}{lrrrc}
    \toprule
    \textbf{Variable}
      & \textbf{AUC Era Antigua}
      & \textbf{AUC Era Nueva}
      & \textbf{$\Delta$AUC}
      & \textbf{Sustancial} \\
    \midrule
    Parcial~1      & 0,869 & 0,781 & $-0{,}087$ & Sí \\
    Parcial~2      & 0,925 & 0,900 & $-0{,}025$ & No \\
    Parcial~3      & 0,931 & 0,845 & $-0{,}086$ & Sí \\
    Prom.\ AMs     & 0,737 & 0,731 & $-0{,}006$ & No \\
    \bottomrule
  \end{tabular}
\end{table}

La Figura~\ref{fig:eras_comparacion} ilustra visualmente esta
comparación, tanto para las variables académicas como para los
patrones de engagement entre eras.

\begin{figure}[htbp]
  \centering
  \includegraphics[width=\textwidth]{09_comparacion_eras.png}
  \caption{\textbf{Izquierda:} AUC univariante de variables
           académicas por era pedagógica. La línea discontinua
           roja marca el umbral mínimo del SAT (AUC~$= 0{,}75$).
           \textbf{Derecha:} medianas de engagement por era
           y resultado académico, normalizadas para comparación
           entre variables de diferente escala.}
  \label{fig:eras_comparacion}
\end{figure}

% ............................................................
\subsubsection{Poder predictivo temprano del Parcial~1:
               tabla de riesgo empírico}
\label{sssec:riesgo_p1}
% ............................................................

El primer checkpoint del SAT (CP1, semana~6) opera únicamente
con la información disponible en ese momento, siendo el Parcial~1
el predictor individual más potente. La
Figura~\ref{fig:riesgo_p1} y la Tabla~\ref{tab:riesgo_p1}
cuantifican el riesgo empírico de reprobación por intervalo
de nota de P1 sobre \texttt{df\_activos} ($N = 554$).

\begin{figure}[htbp]
  \centering
  \includegraphics[width=\textwidth]{03_riesgo_p1.png}
  \caption{\textbf{Izquierda:} diagrama de dispersión entre la
           nota del Parcial~1 y la nota final, coloreado por
           resultado académico. Las líneas discontinuas marcan
           el umbral de aprobación (3{,}0). La correlación de
           Spearman se muestra en el recuadro.
           \textbf{Derecha:} tasa de reprobación empírica por
           intervalo de P1. La línea discontinua horizontal
           indica la tasa global de reprobación (25,7\,\%).}
  \label{fig:riesgo_p1}
\end{figure}

\begin{table}[htbp]
  \centering
  \caption{Tabla de riesgo empírico de reprobación por intervalo
           de nota en Parcial~1 (CP1, semana~6).}
  \label{tab:riesgo_p1}
  \begin{tabular}{ccrrrr}
    \toprule
    \textbf{Intervalo P1}
      & \textbf{Nivel}
      & \textbf{$N$}
      & \textbf{P(APROBÓ)}
      & \textbf{P(REPROBÓ)}
      & \textbf{$\bar{x}$ nota final} \\
    \midrule
    $[0{,}0;\;1{,}5)$  & Crítico  &  19 & 15,8\,\% & \textbf{84,2\,\%} & 2,16 \\
    $[1{,}5;\;2{,}5)$  & Alto     & 102 & 43,1\,\% & \textbf{56,9\,\%} & 2,80 \\
    $[2{,}5;\;3{,}0)$  & Moderado & 100 & 68,0\,\% &         32,0\,\%  & 3,10 \\
    $[3{,}0;\;3{,}5)$  & Límite   & 128 & 85,9\,\% &         14,1\,\%  & 3,44 \\
    $[3{,}5;\;4{,}0)$  & Seguro   &  95 & 83,2\,\% &         16,8\,\%  & 3,62 \\
    $[4{,}0;\;5{,}0]$  & Sólido   & 110 & 99,1\,\% &          0,9\,\%  & 4,18 \\
    \bottomrule
  \end{tabular}
\end{table}

Los resultados son contundentes en los extremos: estudiantes con
P1 $< 1{,}5$ tienen una probabilidad empírica del 84,2\,\% de
reprobar, mientras que con P1~$\geq 4{,}0$ esa probabilidad
cae al 0,9\,\%. Sin embargo, el intervalo $[2{,}5;\;3{,}5)$,
que concentra 228 estudiantes (41,2\,\% del total activo),
presenta tasas de reprobación intermedias (14--32\,\%) donde la
nota del parcial tiene menor poder de discriminación. Este es
precisamente el segmento donde el engagement y otros predictores
complementarios tienen mayor valor incremental.

% ............................................................
\subsubsection{Trayectorias académicas P1~$\to$~P2~$\to$~P3}
\label{sssec:trayectorias}
% ............................................................

Más allá del nivel absoluto de P1, la \emph{trayectoria}
de desempeño a lo largo de los tres parciales provee
información adicional sobre el resultado final. Se identifican
cinco patrones cualitativamente distintos
(Figura~\ref{fig:trayectorias} y Tabla~\ref{tab:trayectorias}).

\begin{figure}[htbp]
  \centering
  \includegraphics[width=\textwidth]{04_trayectorias_parciales.png}
  \caption{Trayectorias de desempeño académico entre parciales.
           \textbf{Izquierda:} frecuencia de cada trayectoria
           por resultado académico.
           \textbf{Derecha:} evolución de la nota promedio
           en cada parcial por trayectoria, con el porcentaje
           de aprobación de cada grupo en la leyenda.
           La línea discontinua marca el umbral de 3{,}0.}
  \label{fig:trayectorias}
\end{figure}

\begin{table}[htbp]
  \centering
  \caption{Patrones de trayectoria académica identificados en
           la secuencia P1~$\to$~P2~$\to$~P3. Los porcentajes
           se calculan dentro de cada grupo de trayectoria.}
  \label{tab:trayectorias}
  \begin{tabular}{lrrrr}
    \toprule
    \textbf{Trayectoria} & \textbf{Total}
      & \textbf{APROBÓ} & \textbf{REPROBÓ} & \textbf{\% APR.} \\
    \midrule
    Declive sostenido (P1~$>$~P2~$>$~P3)    &  40 &  22 &  18 & 55,0\,\% \\
    Estable   (P1~$\approx$~P2~$\approx$~P3)& 306 & 225 &  81 & 73,5\,\% \\
    Mejora sostenida (P1~$<$~P2~$<$~P3)     &  22 &  20 &   2 & 90,9\,\% \\
    Pico intermedio (P1~$<$~P2~$>$~P3)      & 129 &  99 &  30 & 76,7\,\% \\
    Recuperación post-P1 (P1~$<$~P2)        &  57 &  47 &  10 & 82,5\,\% \\
    \bottomrule
  \end{tabular}
\end{table}

El grupo de \textbf{declive sostenido} concentra la mayor tasa
de reprobación (45,0\,\%) y es identificable a partir de la
comparación P1~$\to$~P2 disponible en CP2. La
\textbf{mejora sostenida} es el patrón con mayor probabilidad
de aprobación (90,9\,\%), lo que indica que la trayectoria
ascendente es una señal de éxito más fiable que el nivel
absoluto de P1.

Más relevante aún para el diseño del SAT: 115 de los 221
estudiantes con P1~$< 3{,}0$ terminaron aprobando (52,0\,\%),
con nota media de Parcial~2 $= 3{,}46$ y nota final media
$= 3{,}38$. Este fenómeno de \textit{recuperación académica}
demuestra que un clasificador basado únicamente en P1 tendría
una tasa de falsas alarmas del 52\,\% entre los estudiantes de
riesgo moderado, motivando explícitamente la revisión del estado
de alerta en el segundo checkpoint (CP2, semana~11).

% ------------------------------------------------------------
\subsection{Análisis de la Actividad en Plataforma
            (Engagement)}
\label{subsec:eng}
% ------------------------------------------------------------

El análisis de engagement se realiza sobre
\texttt{df\_activos\_engagement} ($N = 502$): 373 aprobados y
129 reprobados. La Figura~\ref{fig:composicion_modelado}
muestra la composición de este subconjunto.

\begin{figure}[htbp]
  \centering
  \includegraphics[width=0.92\textwidth]{12_composicion_modelado.png}
  \caption{Composición del subconjunto de modelado
           (\texttt{df\_activos\_engagement}, $N = 502$).
           \textbf{Izquierda:} distribución de clases. El
           desbalance 74,3\,\%~/ 25,7\,\% motiva el uso de
           técnicas de balanceo (SMOTE, SMOTEENN).
           \textbf{Derecha:} descomposición por semestre, con
           el porcentaje de reprobados sobre cada barra.}
  \label{fig:composicion_modelado}
\end{figure}

% ............................................................
\subsubsection{Descripción de las métricas de engagement}
\label{sssec:vars_eng}
% ............................................................

Bloque Neón registra el comportamiento de estudio en dos
dimensiones complementarias:

\begin{itemize}
  \item \textbf{Amplitud de estudio (cobertura):}
    módulos únicos visitados (\texttt{parcial\_X\_modulos\_unicos}),
    temas únicos consultados (\texttt{parcial\_X\_temas\_unicos})
    y visitas totales al contenido (\texttt{parcial\_X\_visitas}).

  \item \textbf{Profundidad de estudio (intensidad):}
    tiempo promedio por visita
    (\texttt{parcial\_X\_tiempo\_por\_visita}) y visitas promedio
    por tema (\texttt{parcial\_X\_visitas\_por\_tema}).
\end{itemize}

Además, se dispone del tiempo acumulado hasta cada checkpoint
(\texttt{engagement\_hasta\_p1},
\texttt{engagement\_hasta\_p2}) y del total del semestre
(\texttt{total\_tiempo\_hrs}).

% ............................................................
\subsubsection{Diferencias estadísticas entre aprobados y
               reprobados: prueba de Mann-Whitney}
\label{sssec:mann_eng}
% ............................................................

La Figura~\ref{fig:engagement_boxplots} presenta las
distribuciones de las principales variables de engagement
disponibles en CP1 por semestre y resultado, y la
Tabla~\ref{tab:mw_engagement} reporta los estadísticos
completos de las pruebas Mann-Whitney.

\begin{figure}[htbp]
  \centering
  \includegraphics[width=\textwidth]{07_engagement_boxplots.png}
  \caption{Distribución de las cuatro principales variables de
           engagement en CP1 por semestre y resultado académico.
           El sombreado naranja indica la era antigua y el verde
           la nueva. En cada subgráfico se reportan el $p$-valor
           de Mann-Whitney, el AUC univariante y el coeficiente
           de rango biserial $r_b$.}
  \label{fig:engagement_boxplots}
\end{figure}

\begin{table}[htbp]
  \centering
  \caption{Prueba Mann-Whitney para variables de engagement
           disponibles en CP1. $\tilde{x}$: mediana; $r_b$:
           coeficiente de rango biserial \cite{Kerby2014}.
           Significancia: *** $p < 0{,}001$; * $p < 0{,}05$.}
  \label{tab:mw_engagement}
  \begin{tabular}{lrrrrc}
    \toprule
    \textbf{Variable}
      & \textbf{$\tilde{x}$ APR.}
      & \textbf{$\tilde{x}$ REP.}
      & \textbf{AUC}
      & \textbf{$r_b$}
      & \textbf{$p$-valor} \\
    \midrule
    Tiempo hasta P1 (h)
      & 8,27  & 5,71  & 0,602 & 0,204
      & $2{,}8\times10^{-4}$ *** \\
    Módulos únicos P1
      & 7,00  & 6,00  & 0,601 & 0,203
      & $2{,}7\times10^{-4}$ *** \\
    Visitas contenido P1
      & 37,00 & 32,00 & 0,575 & 0,150
      & $1{,}1\times10^{-2}$ *   \\
    Temas únicos P1
      & 12,00 & 11,00 & 0,567 & 0,134
      & $2{,}3\times10^{-2}$ *   \\
    Visitas por tema P1
      & 3,20  & 2,83  & 0,565 & 0,130
      & $2{,}8\times10^{-2}$ *   \\
    Minutos por visita P1
      & 12,26 & 11,07 & 0,562 & 0,125
      & $3{,}5\times10^{-2}$ *   \\
    \bottomrule
  \end{tabular}
\end{table}

Todas las variables de engagement en CP1 presentan diferencias
estadísticamente significativas ($p < 0{,}05$), pero los
tamaños del efecto son \textbf{pequeños}
($r_b \in [0{,}12;\;0{,}20]$), muy por debajo de los observados
para las variables académicas (Tabla~\ref{tab:mw_academicas}).

\begin{table}[htbp]
  \centering
  \caption{Prueba Mann-Whitney para variables académicas.
           Todas presentan tamaños del efecto grandes
           ($d_{\text{Cohen}} \geq 0{,}91$)
           con significancia $p < 10^{-15}$.}
  \label{tab:mw_academicas}
  \begin{tabular}{lrrrrrr}
    \toprule
    \textbf{Variable}
      & \textbf{$\tilde{x}$ APR.}
      & \textbf{$\tilde{x}$ REP.}
      & \textbf{AUC}
      & \textbf{$r_b$}
      & \textbf{$d_{\text{Cohen}}$}
      & \textbf{$p$-valor} \\
    \midrule
    Parcial~1       & 3,42 & 2,43 & 0,810 & 0,620 & 1,24
      & $1{,}8\times10^{-28}$ \\
    Parcial~2       & 3,80 & 2,33 & 0,861 & 0,722 & 1,56
      & $6{,}7\times10^{-38}$ \\
    Parcial~3       & 3,48 & 1,70 & 0,882 & 0,763 & 1,66
      & $4{,}7\times10^{-42}$ \\
    Prom.\ AMs      & 4,11 & 3,12 & 0,728 & 0,456 & 0,91
      & $2{,}9\times10^{-16}$ \\
    Prom.\ Quizzes  & 4,00 & 2,91 & 0,821 & 0,642 & 1,36
      & $1{,}0\times10^{-20}$ \\
    \bottomrule
  \end{tabular}
\end{table}

% ............................................................
\subsubsection{Poder discriminativo individual: ranking de AUC}
\label{sssec:auc_ranking}
% ............................................................

La Figura~\ref{fig:auc_univariante} resume el poder discriminativo
individual de cada variable disponible en CP1 mediante su AUC
univariante, ordenado de mayor a menor.

\begin{figure}[htbp]
  \centering
  \includegraphics[width=0.92\textwidth]{06_auc_univariante_cp1.png}
  \caption{AUC univariante de cada predictor candidato para CP1.
           Las variables académicas (azul) dominan ampliamente;
           las de engagement (verde) se ubican en la zona
           $[0{,}56;\;0{,}60]$, superando la clasificación
           aleatoria pero por debajo del umbral mínimo del SAT
           (0{,}75, línea verde discontinua). Prom.\ Quizzes no
           está disponible en CP1 para todos los semestres.}
  \label{fig:auc_univariante}
\end{figure}

La brecha entre el AUC de Parcial~1 (0,810) y el del mejor
predictor de engagement (Tiempo hasta P1: 0,602) es de 20,8
puntos porcentuales. Esto confirma que el engagement es un
predictor \emph{complementario}, no un sustituto de las
calificaciones. Su valor radica en los casos marginales donde
la nota de P1 tiene menor poder de discriminación
(zona $[2{,}5;\;3{,}5)$ de la tabla de riesgo).

% ............................................................
\subsubsection{Valor incremental del engagement y correlación
               parcial}
\label{sssec:eng_incremental}
% ............................................................

Para cuantificar el aporte del engagement \emph{más allá} de lo
que ya captura Parcial~1, se calcula la correlación de Spearman
entre el tiempo de engagement hasta P1 y la nota final,
controlando por la nota de P1. El resultado es
$\rho_{\text{parcial}} = 0{,}166$ ($p = 1{,}9 \times 10^{-4}$,
$N = 502$), indicando que el engagement explica un 2--3\,\% de
varianza adicional en el resultado. Aunque modesto en términos
absolutos, este aporte es estadísticamente robusto y
operacionalmente relevante en la zona de incertidumbre del
clasificador.

% ............................................................
\subsubsection{Patrones de acumulación de engagement entre
               checkpoints}
\label{sssec:eng_evolucion}
% ............................................................

La Figura~\ref{fig:engagement_evolucion} explora cómo se acumula
el engagement entre los dos checkpoints del SAT.

\begin{figure}[htbp]
  \centering
  \includegraphics[width=\textwidth]{10_evolucion_engagement.png}
  \caption{\textbf{Izquierda:} correlación entre el tiempo de
           engagement hasta P1 y hasta P2 por resultado académico
           ($\rho_{\text{Spearman}}$ reportado). La diagonal
           discontinua indica el punto de no-incremento entre
           checkpoints.
           \textbf{Derecha:} distribución del incremento de
           engagement en cada período por resultado (violines).
           Los $p$-valores corresponden a pruebas Mann-Whitney
           entre estados.}
  \label{fig:engagement_evolucion}
\end{figure}

La correlación entre el engagement hasta P1 y hasta P2 es alta
para ambos grupos (aprobados: $\rho \approx 0{,}86$;
reprobados: similar), lo que indica \textbf{persistencia de los
hábitos de estudio}: los estudiantes que estudian más en el primer
período tienden a hacerlo también en el segundo. Los aprobados
acumulan en el período P1~$\to$~P2 una mediana de 6,80 horas
adicionales frente a 4,85 de los reprobados; esta diferencia
se amplía hacia el final del semestre (5,29~h vs.~3,28~h en
el período P2~$\to$~fin). Este patrón sugiere que los estudiantes
en riesgo no solo estudian menos, sino que además reducen
progresivamente su dedicación conforme avanza el semestre.

% ------------------------------------------------------------
\subsection{Análisis de Correlaciones y Multicolinealidad}
\label{subsec:corr}
% ------------------------------------------------------------

% ............................................................
\subsubsection{Correlaciones entre variables de engagement}
\label{sssec:corr_eng}
% ............................................................

La Figura~\ref{fig:corr_engagement} presenta la matriz de
correlaciones de Spearman entre las variables de engagement y
la dispersión entre tiempo hasta P1 y nota del Parcial~1.

\begin{figure}[htbp]
  \centering
  \includegraphics[width=\textwidth]{08_correlaciones_engagement.png}
  \caption{\textbf{Izquierda:} matriz de correlaciones de Spearman
           entre variables de engagement (triángulo inferior).
           Las celdas en rojo oscuro indican correlaciones
           altas ($r > 0{,}80$).
           \textbf{Derecha:} dispersión entre tiempo de
           engagement hasta P1 y nota del Parcial~1 por resultado
           académico, con línea de regresión y correlación de
           Spearman.}
  \label{fig:corr_engagement}
\end{figure}

Las correlaciones intra-período son altas: los pares más
problemáticos son tiempo hasta P1 $\leftrightarrow$ tiempo
hasta P2 ($r = 0{,}864$), temas~P1 $\leftrightarrow$
módulos~P1 ($r = 0{,}861$), tiempo~P2 $\leftrightarrow$
tiempo total ($r = 0{,}852$) y temas~total $\leftrightarrow$
módulos~total ($r = 0{,}941$). Esta multicolinealidad no
impide el modelado pero requiere regularización en los
modelos lineales y justifica la selección de subconjuntos
representativos de cada dimensión de engagement en lugar
de incluir todas las variables simultáneamente.

\subsubsection{Factor de inflación de varianza (VIF)}
\label{sssec:vif}

El análisis VIF sobre el conjunto de predictores candidatos
revela que las variables académicas presentan los valores más
altos (Parcial~1: VIF~$= 13{,}7$; AM~2: VIF~$= 10{,}4$;
AM~1: VIF~$= 9{,}4$), mientras que las de engagement son más
moderadas (módulos únicos P1: VIF~$= 7{,}2$;
tiempo hasta P1: VIF~$= 4{,}0$;
tiempo por visita P1: VIF~$= 3{,}9$). Ningún valor supera el
umbral de 15 que produce inestabilidad severa, pero la magnitud
de algunos VIF justifica el uso de regularización L1/L2 en los
modelos de regresión logística (ver Sección~\ref{sec:metodologia}).

% ------------------------------------------------------------
\subsection{Síntesis: Decisiones de Diseño del Dataset}
\label{subsec:sintesis_datos}
% ------------------------------------------------------------

La Tabla~\ref{tab:resumen_decisiones} sintetiza las decisiones
de preprocesamiento tomadas a partir del análisis exploratorio.

\begin{table}[htbp]
  \centering
  \caption{Resumen de decisiones de diseño del dataset con
           justificación empírica e impacto cuantificado.}
  \label{tab:resumen_decisiones}
  \small
  \begin{tabular}{p{3.5cm}p{5.0cm}p{4.8cm}}
    \toprule
    \textbf{Decisión} & \textbf{Justificación} & \textbf{Impacto} \\
    \midrule
    Excluir RETIRADO y SIN~NOTA
      & Etiqueta no observable; sesgo de retiro endógeno;
        heterogeneidad semántica
      & $N$: 764~$\to$~554 \\[3pt]
    Excluir estudiantes sin engagement
      & Cero ambiguo (75\,\% en semestre de arranque);
        $\Delta F_1 < 0{,}02$ al excluirlos
      & $N$: 554~$\to$~502 \\[3pt]
    Excluir variables de actividades evaluativas
      & AUC~$< 0{,}50$; $p > 0{,}05$ (sin discriminación)
      & $-3$ predictores candidatos \\[3pt]
    Incluir \texttt{era\_encoded}
      & $\Delta$AUC(P1) $= -0{,}087$;
        engagement mediano cae 60\,\% en era nueva
      & Control en todos los modelos \\[3pt]
    Incluir \texttt{modalidad\_am}
      & AM-3 grupal vs.\ individual:
        $p = 1{,}4\times10^{-9}$, $d = 0{,}40$
      & Control en modelos con AMs \\[3pt]
    Incluir engagement a pesar de efecto pequeño
      & $\rho_{\text{parcial}} = 0{,}166$, $p < 0{,}001$;
        valor incremental en zona marginal de P1
      & $+6$ predictores en CP1 \\[3pt]
    Aplicar balanceo SMOTE / SMOTEENN
      & 74,3\,\% APROBÓ vs.\ 25,7\,\% REPROBÓ;
        Recall sin balanceo $\approx 0{,}55$
      & Solo en training fold; nunca en evaluación \\
    \bottomrule
  \end{tabular}
\end{table}

El dataset de modelado resultante (\texttt{df\_activos\_engagement},
$N = 502$) presenta tres características estructurales que
condicionan el diseño metodológico de la siguiente sección:
(i)~\textbf{desbalance de clases} (74,3\,\% vs.\ 25,7\,\%),
que exige técnicas de remuestreo para evitar que los modelos
sesguen la predicción hacia la clase mayoritaria;
(ii)~\textbf{estructura longitudinal multisemestral}, que
invalida la validación cruzada aleatoria estándar y exige el
esquema LOSO; y (iii)~\textbf{heterogeneidad de era pedagógica},
que requiere variables de control y análisis estratificado.
Estos tres aspectos se abordan en detalle en la sección de
metodología.

% ============================================================
%  REFERENCIAS
% ============================================================

\begin{thebibliography}{9}

\bibitem{Lopez2021}
J.~López-Zambrano, J.~A. Lara~Torralbo, y C.~Romero,
``Early prediction of student learning performance through data
mining: A systematic review,''
\textit{Psicothema}, vol.~33, no.~3, pp.~456--465, 2021.

\bibitem{Tinto1993}
V.~Tinto,
\textit{Leaving College: Rethinking the Causes and Cures of
Student Attrition}, 2.ª~ed.
Chicago, IL, USA: University of Chicago Press, 1993.

\bibitem{Kerby2014}
D.~S. Kerby,
``The simple difference formula: An approach to teaching
nonparametric correlation,''
\textit{Compr. Psychol.}, vol.~3, p.~11.IT.3.4.1, 2014.

\bibitem{He2009}
H.~He y E.~A. Garcia,
``Learning from imbalanced data,''
\textit{IEEE Trans. Knowl. Data Eng.}, vol.~21, no.~9,
pp.~1263--1284, sep. 2009.

\end{thebibliography}

\end{document}
 y eliminar el
%  preámbulo y los comandos \begin{document}/\end{document}.
%  Cambiar [english] a [spanish] con texlive-lang-spanish.
%  Las figuras se cargan desde la carpeta figures/ relativa.
% ============================================================

\documentclass[12pt, letterpaper]{article}

\usepackage[utf8]{inputenc}
\usepackage[T1]{fontenc}
\usepackage[english]{babel}
\usepackage{amsmath, amssymb}
\usepackage{geometry}
\usepackage{setspace}
\usepackage{booktabs}
\usepackage{multirow}
\usepackage{array}
\usepackage{graphicx}
\usepackage{subcaption}
\usepackage{xcolor}
\usepackage{cite}
\usepackage[hidelinks]{hyperref}
\usepackage{url}
\usepackage{rotating}

\geometry{top=2.5cm, bottom=2.5cm, left=3cm, right=2.5cm}
\onehalfspacing
\renewcommand{\refname}{Referencias}

\graphicspath{{figures/}}

% ============================================================
\begin{document}
% ============================================================

\section{Descripción y Exploración de Datos}
\label{sec:exploracion}

Este capítulo caracteriza el conjunto de datos empleado en el SAT,
documenta las decisiones de preprocesamiento adoptadas y presenta
los hallazgos del análisis exploratorio que fundamentan el diseño
del sistema de modelado. La presentación se organiza en seis
bloques: (i)~fuente y composición del dataset;
(ii)~definición y justificación de los subconjuntos de trabajo;
(iii)~calidad y estructura de los datos;
(iv)~análisis estadístico de las notas académicas;
(v)~análisis de la actividad en plataforma; y
(vi)~análisis de correlaciones y síntesis de decisiones de diseño.

% ------------------------------------------------------------
\subsection{Fuente y Composición del Conjunto de Datos}
\label{subsec:fuente}
% ------------------------------------------------------------

Los datos provienen del curso IIND-2104 \textit{Modelos
Probabilísticos} de la Universidad de los Andes, dictado entre
los semestres 202320 y 202520 —cinco cohortes consecutivas
cubiertas entre 2023 y 2025. La información se integra a partir
de dos fuentes complementarias:

\begin{itemize}
  \item \textbf{Sistema de gestión académica institucional:}
    calificaciones oficiales de los tres parciales y de las
    evaluaciones continuas (actividades modulares, quizzes,
    talleres y proyectos, según el formato de cada semestre).

  \item \textbf{Plataforma Bloque Neón:} sistema de gestión del
    aprendizaje (LMS) que registra, por sesión y por módulo
    de contenido, el tiempo de estudio, los temas consultados,
    los módulos visitados y la frecuencia de acceso de cada
    estudiante.
\end{itemize}

El dataset crudo consolidado cuenta con \textbf{764~registros y
77~variables}, con cuatro posibles resultados académicos finales:
\texttt{APROBÓ} (413 estudiantes, 54,1\,\%),
\texttt{REPROBÓ} (141, 18,5\,\%),
\texttt{RETIRADO} (164, 21,5\,\%) y
\texttt{SIN~NOTA} (46, 6,0\,\%).

La Figura~\ref{fig:composicion} ilustra cómo se distribuyen
estos resultados por semestre en términos absolutos y
porcentuales, y la Tabla~\ref{tab:composicion_semestral} detalla
las frecuencias exactas.

\begin{figure}[htbp]
  \centering
  \includegraphics[width=\textwidth]{01_composicion_dataset.png}
  \caption{Composición del dataset por semestre y resultado
           académico. \textbf{Izquierda:} distribución absoluta.
           \textbf{Derecha:} distribución porcentual con
           etiquetas de proporción. El sombreado naranja
           corresponde a la era pedagógica antigua (202320--202410)
           y el verde a la nueva (202420--202520).}
  \label{fig:composicion}
\end{figure}

\begin{table}[htbp]
  \centering
  \caption{Composición por semestre, era pedagógica y resultado
           académico final.}
  \label{tab:composicion_semestral}
  \small
  \begin{tabular}{lccrrrrrrr}
    \toprule
    \textbf{Semestre} & \textbf{Era} & \textbf{Total}
      & \multicolumn{2}{c}{\textbf{APROBÓ}}
      & \multicolumn{2}{c}{\textbf{REPROBÓ}}
      & \multicolumn{2}{c}{\textbf{RETIRADO}}
      & \textbf{SIN} \\
    \cmidrule(lr){4-5}\cmidrule(lr){6-7}\cmidrule(lr){8-9}
    & & & $n$ & \% & $n$ & \% & $n$ & \% & \textbf{NOTA} \\
    \midrule
    202320 & Antigua & 170 & 116 & 68,2 &  27 & 15,9 &  27 & 15,9 & 0  \\
    202410 & Antigua & 150 &  68 & 45,3 &  14 &  9,3 &  30 & 20,0 & 38 \\
    202420 & Nueva   & 164 &  90 & 54,9 &  33 & 20,1 &  40 & 24,4 & 1  \\
    202510 & Nueva   & 140 &  70 & 50,0 &  47 & 33,6 &  21 & 15,0 & 2  \\
    202520 & Nueva   & 140 &  69 & 49,3 &  20 & 14,3 &  46 & 32,9 & 5  \\
    \midrule
    \textbf{Total} & — & \textbf{764}
      & \textbf{413} & \textbf{54,1}
      & \textbf{141} & \textbf{18,5}
      & \textbf{164} & \textbf{21,5}
      & \textbf{46} \\
    \bottomrule
  \end{tabular}
\end{table}

Dos fenómenos merecen atención especial. Primero, la tasa de
reprobación en 202510 (33,6\,\%) es la más alta del período y
más del doble de la de 202410 (9,3\,\%), lo que sugiere
variabilidad intercohorte que el modelo debe generalizar. Segundo,
el semestre 202520 registra la mayor tasa de retiro (32,9\,\%),
fenómeno atípico que podría obedecer a cambios en la política
institucional de permanencia o a condiciones externas no
capturadas en los datos.

% ............................................................
\subsubsection{Dos eras pedagógicas}
\label{sssec:eras}
% ............................................................

Un rasgo estructural del dataset es la coexistencia de dos
\textbf{formatos de evaluación continua}:

\begin{itemize}
  \item \textbf{Era Antigua} (202320 y 202410; $n = 320$ estudiantes
    activos): las evaluaciones continuas consisten en talleres
    individuales y proyectos por fases grupales.

  \item \textbf{Era Nueva} (202420, 202510 y 202520; $n = 444$
    estudiantes activos): las evaluaciones continuas son
    actividades modulares (AMs) y quizzes individuales.
\end{itemize}

Dentro de la era nueva existe una distinción adicional: en
202420 las AMs se realizaron \textbf{en grupo}, mientras que
a partir de 202510 pasaron a ser \textbf{individuales}. Este
cambio produce diferencias estadísticamente significativas en la
distribución de AM-3 ($U$-test, $p = 1{,}39 \times 10^{-9}$,
$d_{\text{Cohen}} = 0{,}40$, AUC$_{\text{grupo}} = 0{,}369$
vs.\ AUC$_{\text{ind.}} = 0{,}246$), lo que justifica incluir la
variable \texttt{modalidad\_am} como control en los modelos
entrenados con datos del formato nuevo.

% ------------------------------------------------------------
\subsection{Definición y Justificación de los Subconjuntos
            de Trabajo}
\label{subsec:subsets}
% ------------------------------------------------------------

El dataset bruto requiere dos depuraciones sucesivas antes de
ser apto para el modelado de clasificación.

% ............................................................
\subsubsection{Exclusión de estudiantes retirados y sin
               calificación}
\label{sssec:excl_retirados}
% ............................................................

Los 164 estudiantes con resultado \texttt{RETIRADO} y los
46 con \texttt{SIN~NOTA} se excluyen por tres razones:

\begin{enumerate}
  \item \textbf{Ausencia de etiqueta objetivo observable.}
    El SAT modela $P(\text{REPROBÓ} \mid \mathbf{x}_i)$, la
    probabilidad de que un estudiante \emph{finalice} el semestre
    reprobando. Los estudiantes retirados abandonaron el curso
    antes de producirse el resultado; su etiqueta final es
    desconocida y no puede inferirse sin supuestos fuertes.

  \item \textbf{Sesgo de retiro endógeno.}
    La decisión de retirarse está correlacionada con el
    desempeño anticipado: los estudiantes que perciben que van
    a reprobar se retiran en mayor proporción \cite{Tinto1993}.
    Incluirlos en la clase \texttt{REPROBÓ} introduciría sesgo
    de selección en las variables de engagement: un estudiante
    retirado en la semana~4 tiene engagement mínimo no por falta
    de habilidad sino por abandono temprano, contaminando el
    predictor más utilizado.

  \item \textbf{Heterogeneidad semántica.}
    La predicción del retiro como fenómeno independiente
    requeriría modelos multi-clase y supuestos adicionales
    sobre su equivalencia con la reprobación, lo cual está
    fuera del alcance del presente trabajo.
\end{enumerate}

El subconjunto resultante es \textbf{\texttt{df\_activos}}
($N = 554$; APROBÓ: 413, REPROBÓ: 141).

% ............................................................
\subsubsection{Exclusión adicional de estudiantes sin actividad
               en plataforma}
\label{sssec:excl_engagement}
% ............................................................

Del subconjunto \texttt{df\_activos}, 52 estudiantes (9,4\,\%)
presentan \textbf{engagement total nulo} en Bloque Neón durante
todo el semestre (i.e., \texttt{engagement\_hasta\_p1} $= 0$).
Estos casos se excluyen del subconjunto de modelado por las
razones siguientes:

\begin{enumerate}
  \item \textbf{Ambigüedad del cero.}
    Existen dos interpretaciones mutuamente excluyentes:
    (a)~el estudiante genuinamente no utilizó la plataforma,
    o (b)~ocurrió un error en el pipeline de integración de
    datos. Ambas hipótesis son plausibles: 39 de los 52 casos
    (75\,\%) se concentran en el semestre 202320, el primero
    en incorporar Bloque Neón, cuando el sistema estaba aún
    en fase de implementación. No es posible distinguir entre
    causas con la información disponible.

  \item \textbf{Riesgo de imputación espuria.}
    Imputar cero como valor real de engagement asumiría que
    ``no usar la plataforma'' es una señal homogénea entre
    semestres, cuando los patrones de no-uso varían
    considerablemente entre cohortes.

  \item \textbf{Impacto estadístico marginal confirmado.}
    Se realizó un análisis de sensibilidad: los modelos
    entrenados sobre \texttt{df\_activos} completo ($N = 554$)
    y sobre el subconjunto filtrado ($N = 502$) arrojan
    diferencias en F1 menores a $\Delta F_1 < 0{,}02$,
    confirmando que la exclusión no altera las conclusiones.
\end{enumerate}

El subconjunto de modelado final es
\textbf{\texttt{df\_activos\_engagement}}
($N = 502$; APROBÓ:~373, REPROBÓ:~129).
La Tabla~\ref{tab:subsets} resume los tres subconjuntos
empleados en el análisis.

\begin{table}[htbp]
  \centering
  \caption{Subconjuntos de trabajo. \texttt{df\_activos\_engagement}
           es el conjunto empleado para el entrenamiento y
           evaluación de todos los modelos del SAT.}
  \label{tab:subsets}
  \begin{tabular}{lrrrl}
    \toprule
    \textbf{Subconjunto} & \textbf{$N$}
      & \textbf{APROBÓ} & \textbf{REPROBÓ}
      & \textbf{Uso principal} \\
    \midrule
    \texttt{df\_all}
      & 764 & 413 & 141
      & Tasas globales, análisis de retiro \\
    \texttt{df\_activos}
      & 554 & 413 & 141
      & Análisis de notas y distribuciones \\
    \texttt{df\_activos\_engagement}
      & \textbf{502} & \textbf{373} & \textbf{129}
      & \textbf{Entrenamiento y evaluación de modelos} \\
    \bottomrule
  \end{tabular}
\end{table}

% ------------------------------------------------------------
\subsection{Calidad y Estructura de los Datos}
\label{subsec:calidad}
% ------------------------------------------------------------

% ............................................................
\subsubsection{Nulos estructurales versus problemáticos}
\label{sssec:nulos}
% ............................................................

El dataset presenta un volumen considerable de valores faltantes,
en su mayoría atribuibles al cambio de formato pedagógico.
Se distinguen dos categorías.

\paragraph{Nulos estructurales (esperados).}
Son variables que solo existen en una era y, por tanto, están
ausentes de forma legítima en la otra:
\texttt{quiz\_4/5/6} (77,8\,\% de nulos, solo en algunos
semestres del formato nuevo),
\texttt{taller\_1/2/3} (59,4\,\%, exclusivo de la era antigua),
\texttt{fase\_1/2} (59,4\,\%, proyectos grupales de la era antigua),
y \texttt{am\_4/5/6} (56,7\,\%, AMs opcionales no presentes en
todos los semestres). Estos nulos no representan problemas de
calidad sino de diseño del instrumento de evaluación.

\paragraph{Nulos problemáticos (casos especiales).}
Son faltantes en variables que deberían estar presentes para
todos los estudiantes activos:
\begin{itemize}
  \item \textbf{Parcial~1 $= 0$ ($n = 1$):} un único estudiante
    registra nota cero. No es posible distinguir ausencia de nota
    real. Se retiene con variable indicadora
    \texttt{no\_presento\_p1}.
  \item \textbf{Parcial~3 $= 0$ ($n = 15$; 13 reprobaron):}
    concentrado en estudiantes encaminados a reprobar, consistente
    con ausentismo por desmotivación al final del semestre.
    Se retienen con su valor real.
  \item \textbf{Notas superiores a 5{,}0:} presentes por
    bonificaciones otorgadas por docentes. Son valores legítimos
    y se retienen sin transformación.
\end{itemize}

% ............................................................
\subsubsection{Variables de plataforma sin poder discriminativo}
\label{sssec:excl_vars}
% ............................................................

La plataforma Bloque Neón registra tanto métricas de
\textit{estudio de contenidos} como métricas de
\textit{actividades evaluativas en línea}. El análisis estadístico
univariante sobre estas últimas —\texttt{actividades\_n\_realizadas},
\texttt{actividades\_tiempo\_total} y
\texttt{actividades\_visitas\_total}— revela ausencia de
poder discriminativo entre aprobados y reprobados
(Tabla~\ref{tab:excl_vars}), lo que se ilustra también en la
Figura~\ref{fig:vars_excluidas}.

\begin{table}[htbp]
  \centering
  \caption{Variables de actividades evaluativas de la plataforma.
           AUC~$\approx 0{,}5$ y $p > 0{,}05$ indican que no
           discriminan mejor que la clasificación aleatoria.}
  \label{tab:excl_vars}
  \begin{tabular}{lrrrr}
    \toprule
    \textbf{Variable}
      & \textbf{Med.\ APROBÓ}
      & \textbf{Med.\ REPROBÓ}
      & \textbf{AUC}
      & \textbf{$p$-valor (MW)} \\
    \midrule
    \texttt{actividades\_n\_realizadas}   & 1,00 & 1,00 & 0,473 & 0,830 \\
    \texttt{actividades\_tiempo\_total}   & 1,15 & 1,10 & 0,498 & 0,532 \\
    \texttt{actividades\_visitas\_total}  & 9,00 & 11,00 & 0,465 & 0,886 \\
    \bottomrule
  \end{tabular}
\end{table}

\begin{figure}[htbp]
  \centering
  \includegraphics[width=\textwidth]{11_variables_excluidas.png}
  \caption{Distribuciones de las tres variables de actividades
           evaluativas para aprobados (azul) y reprobados (rojo).
           Los valores de $p > 0{,}05$ y AUC~$< 0{,}5$
           confirman la ausencia de poder discriminativo,
           justificando su exclusión del conjunto de predictores.}
  \label{fig:vars_excluidas}
\end{figure}

Nótese que, para \texttt{actividades\_visitas\_total}, la mediana
de los reprobados (11) supera la de los aprobados (9), con
AUC~$= 0{,}465$ —peor que la clasificación aleatoria—, lo que
sugiere que los estudiantes con mayor número de visitas a
actividades no necesariamente obtienen mejores resultados. Las
tres variables son excluidas del conjunto de predictores para
evitar introducir ruido en los modelos.

% ------------------------------------------------------------
\subsection{Análisis Estadístico de las Notas Académicas}
\label{subsec:notas}
% ------------------------------------------------------------

El análisis de notas se realiza sobre \texttt{df\_activos}
($N = 554$) para preservar la distribución real sin los
filtros de engagement.

% ............................................................
\subsubsection{Distribución de la nota final}
\label{sssec:dist_nota}
% ............................................................

La nota final (escala 0--5) presenta una distribución bimodal
con dos concentraciones claramente separadas: una en torno a
2,6 (zona de reprobación) y otra en torno a 3,7 (zona de
aprobación cómoda), como se observa en la
Figura~\ref{fig:dist_nota}.

\begin{figure}[htbp]
  \centering
  \includegraphics[width=\textwidth]{02_distribucion_nota_final.png}
  \caption{Distribución de la nota final.
           \textbf{Izquierda:} histogramas superpuestos con
           curva de densidad KDE por resultado académico.
           La línea discontinua vertical indica el umbral de
           aprobación (3{,}0).
           \textbf{Derecha:} diagramas de caja por semestre y
           resultado; el sombreado distingue las eras pedagógicas.}
  \label{fig:dist_nota}
\end{figure}

Las estadísticas descriptivas detalladas se presentan en la
Tabla~\ref{tab:desc_nota}. La prueba de Mann-Whitney confirma
la separación altamente significativa entre grupos
($U = 58\,233{,}0$, $p = 9{,}97 \times 10^{-71}$,
$d_{\text{Cohen}} = 2{,}41$). Esta magnitud del efecto implica
que la distribución de notas de los aprobados y la de los
reprobados se solapan apenas en una pequeña región central
(alrededor de $[2{,}8;\;3{,}2]$), que es precisamente la zona
donde los modelos de clasificación concentran su incertidumbre.

\begin{table}[htbp]
  \centering
  \caption{Estadísticas descriptivas de la nota final por resultado
           y era pedagógica. $\bar{x}$: media; $\tilde{x}$: mediana;
           $s$: desviación estándar.}
  \label{tab:desc_nota}
  \begin{tabular}{lrrrr}
    \toprule
    \textbf{Grupo} & \textbf{$N$}
      & \textbf{$\bar{x}$} & \textbf{$\tilde{x}$} & \textbf{$s$} \\
    \midrule
    APROBÓ (global)                & 413 & 3,724 & 3,620 & 0,527 \\
    REPROBÓ (global)               & 141 & 2,418 & 2,632 & 0,583 \\
    \midrule
    Era Antigua — APROBÓ           & 184 & 3,777 & 3,680 & 0,504 \\
    Era Antigua — REPROBÓ          &  41 & 2,446 & 2,720 & 0,671 \\
    Era Nueva   — APROBÓ           & 229 & 3,682 & 3,530 & 0,541 \\
    Era Nueva   — REPROBÓ          & 100 & 2,407 & 2,605 & 0,545 \\
    \bottomrule
  \end{tabular}
\end{table}

% ............................................................
\subsubsection{Poder predictivo de los parciales:
               análisis por era pedagógica}
\label{sssec:parciales_era}
% ............................................................

La Figura~\ref{fig:parciales_era} compara la distribución de
las tres calificaciones de parciales entre aprobados y reprobados,
estratificada por era pedagógica.

\begin{figure}[htbp]
  \centering
  \includegraphics[width=\textwidth]{05_parciales_era_resultado.png}
  \caption{Distribuciones de Parcial~1, Parcial~2 y Parcial~3
           por resultado académico, estratificadas por era
           pedagógica. Cada subgráfico incluye el $p$-valor
           de la prueba Mann-Whitney, el AUC univariante y el
           coeficiente de rango biserial~$r_b$.}
  \label{fig:parciales_era}
\end{figure}

La Tabla~\ref{tab:auc_era} sintetiza los AUC univariantes por
variable y era. El hallazgo más relevante es la \textbf{reducción
del poder predictivo de P1 en la era nueva}
($\Delta\text{AUC} = -0{,}087$). Esta caída puede atribuirse
al mayor peso relativo de las actividades continuas (quizzes y
AMs) en el nuevo currículo, que diluye la señal del primer
examen. Contrariamente, el Parcial~2 mantiene una capacidad
predictiva alta y estable en ambas eras
($\Delta\text{AUC} = -0{,}025$), lo que refuerza el valor
del segundo checkpoint (CP2) del SAT.

\begin{table}[htbp]
  \centering
  \caption{AUC univariante de las variables académicas por era
           pedagógica. $\Delta$AUC indica la variación entre eras.
           Diferencias $> 5\,\%$ se consideran sustanciales.}
  \label{tab:auc_era}
  \begin{tabular}{lrrrc}
    \toprule
    \textbf{Variable}
      & \textbf{AUC Era Antigua}
      & \textbf{AUC Era Nueva}
      & \textbf{$\Delta$AUC}
      & \textbf{Sustancial} \\
    \midrule
    Parcial~1      & 0,869 & 0,781 & $-0{,}087$ & Sí \\
    Parcial~2      & 0,925 & 0,900 & $-0{,}025$ & No \\
    Parcial~3      & 0,931 & 0,845 & $-0{,}086$ & Sí \\
    Prom.\ AMs     & 0,737 & 0,731 & $-0{,}006$ & No \\
    \bottomrule
  \end{tabular}
\end{table}

La Figura~\ref{fig:eras_comparacion} ilustra visualmente esta
comparación, tanto para las variables académicas como para los
patrones de engagement entre eras.

\begin{figure}[htbp]
  \centering
  \includegraphics[width=\textwidth]{09_comparacion_eras.png}
  \caption{\textbf{Izquierda:} AUC univariante de variables
           académicas por era pedagógica. La línea discontinua
           roja marca el umbral mínimo del SAT (AUC~$= 0{,}75$).
           \textbf{Derecha:} medianas de engagement por era
           y resultado académico, normalizadas para comparación
           entre variables de diferente escala.}
  \label{fig:eras_comparacion}
\end{figure}

% ............................................................
\subsubsection{Poder predictivo temprano del Parcial~1:
               tabla de riesgo empírico}
\label{sssec:riesgo_p1}
% ............................................................

El primer checkpoint del SAT (CP1, semana~6) opera únicamente
con la información disponible en ese momento, siendo el Parcial~1
el predictor individual más potente. La
Figura~\ref{fig:riesgo_p1} y la Tabla~\ref{tab:riesgo_p1}
cuantifican el riesgo empírico de reprobación por intervalo
de nota de P1 sobre \texttt{df\_activos} ($N = 554$).

\begin{figure}[htbp]
  \centering
  \includegraphics[width=\textwidth]{03_riesgo_p1.png}
  \caption{\textbf{Izquierda:} diagrama de dispersión entre la
           nota del Parcial~1 y la nota final, coloreado por
           resultado académico. Las líneas discontinuas marcan
           el umbral de aprobación (3{,}0). La correlación de
           Spearman se muestra en el recuadro.
           \textbf{Derecha:} tasa de reprobación empírica por
           intervalo de P1. La línea discontinua horizontal
           indica la tasa global de reprobación (25,7\,\%).}
  \label{fig:riesgo_p1}
\end{figure}

\begin{table}[htbp]
  \centering
  \caption{Tabla de riesgo empírico de reprobación por intervalo
           de nota en Parcial~1 (CP1, semana~6).}
  \label{tab:riesgo_p1}
  \begin{tabular}{ccrrrr}
    \toprule
    \textbf{Intervalo P1}
      & \textbf{Nivel}
      & \textbf{$N$}
      & \textbf{P(APROBÓ)}
      & \textbf{P(REPROBÓ)}
      & \textbf{$\bar{x}$ nota final} \\
    \midrule
    $[0{,}0;\;1{,}5)$  & Crítico  &  19 & 15,8\,\% & \textbf{84,2\,\%} & 2,16 \\
    $[1{,}5;\;2{,}5)$  & Alto     & 102 & 43,1\,\% & \textbf{56,9\,\%} & 2,80 \\
    $[2{,}5;\;3{,}0)$  & Moderado & 100 & 68,0\,\% &         32,0\,\%  & 3,10 \\
    $[3{,}0;\;3{,}5)$  & Límite   & 128 & 85,9\,\% &         14,1\,\%  & 3,44 \\
    $[3{,}5;\;4{,}0)$  & Seguro   &  95 & 83,2\,\% &         16,8\,\%  & 3,62 \\
    $[4{,}0;\;5{,}0]$  & Sólido   & 110 & 99,1\,\% &          0,9\,\%  & 4,18 \\
    \bottomrule
  \end{tabular}
\end{table}

Los resultados son contundentes en los extremos: estudiantes con
P1 $< 1{,}5$ tienen una probabilidad empírica del 84,2\,\% de
reprobar, mientras que con P1~$\geq 4{,}0$ esa probabilidad
cae al 0,9\,\%. Sin embargo, el intervalo $[2{,}5;\;3{,}5)$,
que concentra 228 estudiantes (41,2\,\% del total activo),
presenta tasas de reprobación intermedias (14--32\,\%) donde la
nota del parcial tiene menor poder de discriminación. Este es
precisamente el segmento donde el engagement y otros predictores
complementarios tienen mayor valor incremental.

% ............................................................
\subsubsection{Trayectorias académicas P1~$\to$~P2~$\to$~P3}
\label{sssec:trayectorias}
% ............................................................

Más allá del nivel absoluto de P1, la \emph{trayectoria}
de desempeño a lo largo de los tres parciales provee
información adicional sobre el resultado final. Se identifican
cinco patrones cualitativamente distintos
(Figura~\ref{fig:trayectorias} y Tabla~\ref{tab:trayectorias}).

\begin{figure}[htbp]
  \centering
  \includegraphics[width=\textwidth]{04_trayectorias_parciales.png}
  \caption{Trayectorias de desempeño académico entre parciales.
           \textbf{Izquierda:} frecuencia de cada trayectoria
           por resultado académico.
           \textbf{Derecha:} evolución de la nota promedio
           en cada parcial por trayectoria, con el porcentaje
           de aprobación de cada grupo en la leyenda.
           La línea discontinua marca el umbral de 3{,}0.}
  \label{fig:trayectorias}
\end{figure}

\begin{table}[htbp]
  \centering
  \caption{Patrones de trayectoria académica identificados en
           la secuencia P1~$\to$~P2~$\to$~P3. Los porcentajes
           se calculan dentro de cada grupo de trayectoria.}
  \label{tab:trayectorias}
  \begin{tabular}{lrrrr}
    \toprule
    \textbf{Trayectoria} & \textbf{Total}
      & \textbf{APROBÓ} & \textbf{REPROBÓ} & \textbf{\% APR.} \\
    \midrule
    Declive sostenido (P1~$>$~P2~$>$~P3)    &  40 &  22 &  18 & 55,0\,\% \\
    Estable   (P1~$\approx$~P2~$\approx$~P3)& 306 & 225 &  81 & 73,5\,\% \\
    Mejora sostenida (P1~$<$~P2~$<$~P3)     &  22 &  20 &   2 & 90,9\,\% \\
    Pico intermedio (P1~$<$~P2~$>$~P3)      & 129 &  99 &  30 & 76,7\,\% \\
    Recuperación post-P1 (P1~$<$~P2)        &  57 &  47 &  10 & 82,5\,\% \\
    \bottomrule
  \end{tabular}
\end{table}

El grupo de \textbf{declive sostenido} concentra la mayor tasa
de reprobación (45,0\,\%) y es identificable a partir de la
comparación P1~$\to$~P2 disponible en CP2. La
\textbf{mejora sostenida} es el patrón con mayor probabilidad
de aprobación (90,9\,\%), lo que indica que la trayectoria
ascendente es una señal de éxito más fiable que el nivel
absoluto de P1.

Más relevante aún para el diseño del SAT: 115 de los 221
estudiantes con P1~$< 3{,}0$ terminaron aprobando (52,0\,\%),
con nota media de Parcial~2 $= 3{,}46$ y nota final media
$= 3{,}38$. Este fenómeno de \textit{recuperación académica}
demuestra que un clasificador basado únicamente en P1 tendría
una tasa de falsas alarmas del 52\,\% entre los estudiantes de
riesgo moderado, motivando explícitamente la revisión del estado
de alerta en el segundo checkpoint (CP2, semana~11).

% ------------------------------------------------------------
\subsection{Análisis de la Actividad en Plataforma
            (Engagement)}
\label{subsec:eng}
% ------------------------------------------------------------

El análisis de engagement se realiza sobre
\texttt{df\_activos\_engagement} ($N = 502$): 373 aprobados y
129 reprobados. La Figura~\ref{fig:composicion_modelado}
muestra la composición de este subconjunto.

\begin{figure}[htbp]
  \centering
  \includegraphics[width=0.92\textwidth]{12_composicion_modelado.png}
  \caption{Composición del subconjunto de modelado
           (\texttt{df\_activos\_engagement}, $N = 502$).
           \textbf{Izquierda:} distribución de clases. El
           desbalance 74,3\,\%~/ 25,7\,\% motiva el uso de
           técnicas de balanceo (SMOTE, SMOTEENN).
           \textbf{Derecha:} descomposición por semestre, con
           el porcentaje de reprobados sobre cada barra.}
  \label{fig:composicion_modelado}
\end{figure}

% ............................................................
\subsubsection{Descripción de las métricas de engagement}
\label{sssec:vars_eng}
% ............................................................

Bloque Neón registra el comportamiento de estudio en dos
dimensiones complementarias:

\begin{itemize}
  \item \textbf{Amplitud de estudio (cobertura):}
    módulos únicos visitados (\texttt{parcial\_X\_modulos\_unicos}),
    temas únicos consultados (\texttt{parcial\_X\_temas\_unicos})
    y visitas totales al contenido (\texttt{parcial\_X\_visitas}).

  \item \textbf{Profundidad de estudio (intensidad):}
    tiempo promedio por visita
    (\texttt{parcial\_X\_tiempo\_por\_visita}) y visitas promedio
    por tema (\texttt{parcial\_X\_visitas\_por\_tema}).
\end{itemize}

Además, se dispone del tiempo acumulado hasta cada checkpoint
(\texttt{engagement\_hasta\_p1},
\texttt{engagement\_hasta\_p2}) y del total del semestre
(\texttt{total\_tiempo\_hrs}).

% ............................................................
\subsubsection{Diferencias estadísticas entre aprobados y
               reprobados: prueba de Mann-Whitney}
\label{sssec:mann_eng}
% ............................................................

La Figura~\ref{fig:engagement_boxplots} presenta las
distribuciones de las principales variables de engagement
disponibles en CP1 por semestre y resultado, y la
Tabla~\ref{tab:mw_engagement} reporta los estadísticos
completos de las pruebas Mann-Whitney.

\begin{figure}[htbp]
  \centering
  \includegraphics[width=\textwidth]{07_engagement_boxplots.png}
  \caption{Distribución de las cuatro principales variables de
           engagement en CP1 por semestre y resultado académico.
           El sombreado naranja indica la era antigua y el verde
           la nueva. En cada subgráfico se reportan el $p$-valor
           de Mann-Whitney, el AUC univariante y el coeficiente
           de rango biserial $r_b$.}
  \label{fig:engagement_boxplots}
\end{figure}

\begin{table}[htbp]
  \centering
  \caption{Prueba Mann-Whitney para variables de engagement
           disponibles en CP1. $\tilde{x}$: mediana; $r_b$:
           coeficiente de rango biserial \cite{Kerby2014}.
           Significancia: *** $p < 0{,}001$; * $p < 0{,}05$.}
  \label{tab:mw_engagement}
  \begin{tabular}{lrrrrc}
    \toprule
    \textbf{Variable}
      & \textbf{$\tilde{x}$ APR.}
      & \textbf{$\tilde{x}$ REP.}
      & \textbf{AUC}
      & \textbf{$r_b$}
      & \textbf{$p$-valor} \\
    \midrule
    Tiempo hasta P1 (h)
      & 8,27  & 5,71  & 0,602 & 0,204
      & $2{,}8\times10^{-4}$ *** \\
    Módulos únicos P1
      & 7,00  & 6,00  & 0,601 & 0,203
      & $2{,}7\times10^{-4}$ *** \\
    Visitas contenido P1
      & 37,00 & 32,00 & 0,575 & 0,150
      & $1{,}1\times10^{-2}$ *   \\
    Temas únicos P1
      & 12,00 & 11,00 & 0,567 & 0,134
      & $2{,}3\times10^{-2}$ *   \\
    Visitas por tema P1
      & 3,20  & 2,83  & 0,565 & 0,130
      & $2{,}8\times10^{-2}$ *   \\
    Minutos por visita P1
      & 12,26 & 11,07 & 0,562 & 0,125
      & $3{,}5\times10^{-2}$ *   \\
    \bottomrule
  \end{tabular}
\end{table}

Todas las variables de engagement en CP1 presentan diferencias
estadísticamente significativas ($p < 0{,}05$), pero los
tamaños del efecto son \textbf{pequeños}
($r_b \in [0{,}12;\;0{,}20]$), muy por debajo de los observados
para las variables académicas (Tabla~\ref{tab:mw_academicas}).

\begin{table}[htbp]
  \centering
  \caption{Prueba Mann-Whitney para variables académicas.
           Todas presentan tamaños del efecto grandes
           ($d_{\text{Cohen}} \geq 0{,}91$)
           con significancia $p < 10^{-15}$.}
  \label{tab:mw_academicas}
  \begin{tabular}{lrrrrrr}
    \toprule
    \textbf{Variable}
      & \textbf{$\tilde{x}$ APR.}
      & \textbf{$\tilde{x}$ REP.}
      & \textbf{AUC}
      & \textbf{$r_b$}
      & \textbf{$d_{\text{Cohen}}$}
      & \textbf{$p$-valor} \\
    \midrule
    Parcial~1       & 3,42 & 2,43 & 0,810 & 0,620 & 1,24
      & $1{,}8\times10^{-28}$ \\
    Parcial~2       & 3,80 & 2,33 & 0,861 & 0,722 & 1,56
      & $6{,}7\times10^{-38}$ \\
    Parcial~3       & 3,48 & 1,70 & 0,882 & 0,763 & 1,66
      & $4{,}7\times10^{-42}$ \\
    Prom.\ AMs      & 4,11 & 3,12 & 0,728 & 0,456 & 0,91
      & $2{,}9\times10^{-16}$ \\
    Prom.\ Quizzes  & 4,00 & 2,91 & 0,821 & 0,642 & 1,36
      & $1{,}0\times10^{-20}$ \\
    \bottomrule
  \end{tabular}
\end{table}

% ............................................................
\subsubsection{Poder discriminativo individual: ranking de AUC}
\label{sssec:auc_ranking}
% ............................................................

La Figura~\ref{fig:auc_univariante} resume el poder discriminativo
individual de cada variable disponible en CP1 mediante su AUC
univariante, ordenado de mayor a menor.

\begin{figure}[htbp]
  \centering
  \includegraphics[width=0.92\textwidth]{06_auc_univariante_cp1.png}
  \caption{AUC univariante de cada predictor candidato para CP1.
           Las variables académicas (azul) dominan ampliamente;
           las de engagement (verde) se ubican en la zona
           $[0{,}56;\;0{,}60]$, superando la clasificación
           aleatoria pero por debajo del umbral mínimo del SAT
           (0{,}75, línea verde discontinua). Prom.\ Quizzes no
           está disponible en CP1 para todos los semestres.}
  \label{fig:auc_univariante}
\end{figure}

La brecha entre el AUC de Parcial~1 (0,810) y el del mejor
predictor de engagement (Tiempo hasta P1: 0,602) es de 20,8
puntos porcentuales. Esto confirma que el engagement es un
predictor \emph{complementario}, no un sustituto de las
calificaciones. Su valor radica en los casos marginales donde
la nota de P1 tiene menor poder de discriminación
(zona $[2{,}5;\;3{,}5)$ de la tabla de riesgo).

% ............................................................
\subsubsection{Valor incremental del engagement y correlación
               parcial}
\label{sssec:eng_incremental}
% ............................................................

Para cuantificar el aporte del engagement \emph{más allá} de lo
que ya captura Parcial~1, se calcula la correlación de Spearman
entre el tiempo de engagement hasta P1 y la nota final,
controlando por la nota de P1. El resultado es
$\rho_{\text{parcial}} = 0{,}166$ ($p = 1{,}9 \times 10^{-4}$,
$N = 502$), indicando que el engagement explica un 2--3\,\% de
varianza adicional en el resultado. Aunque modesto en términos
absolutos, este aporte es estadísticamente robusto y
operacionalmente relevante en la zona de incertidumbre del
clasificador.

% ............................................................
\subsubsection{Patrones de acumulación de engagement entre
               checkpoints}
\label{sssec:eng_evolucion}
% ............................................................

La Figura~\ref{fig:engagement_evolucion} explora cómo se acumula
el engagement entre los dos checkpoints del SAT.

\begin{figure}[htbp]
  \centering
  \includegraphics[width=\textwidth]{10_evolucion_engagement.png}
  \caption{\textbf{Izquierda:} correlación entre el tiempo de
           engagement hasta P1 y hasta P2 por resultado académico
           ($\rho_{\text{Spearman}}$ reportado). La diagonal
           discontinua indica el punto de no-incremento entre
           checkpoints.
           \textbf{Derecha:} distribución del incremento de
           engagement en cada período por resultado (violines).
           Los $p$-valores corresponden a pruebas Mann-Whitney
           entre estados.}
  \label{fig:engagement_evolucion}
\end{figure}

La correlación entre el engagement hasta P1 y hasta P2 es alta
para ambos grupos (aprobados: $\rho \approx 0{,}86$;
reprobados: similar), lo que indica \textbf{persistencia de los
hábitos de estudio}: los estudiantes que estudian más en el primer
período tienden a hacerlo también en el segundo. Los aprobados
acumulan en el período P1~$\to$~P2 una mediana de 6,80 horas
adicionales frente a 4,85 de los reprobados; esta diferencia
se amplía hacia el final del semestre (5,29~h vs.~3,28~h en
el período P2~$\to$~fin). Este patrón sugiere que los estudiantes
en riesgo no solo estudian menos, sino que además reducen
progresivamente su dedicación conforme avanza el semestre.

% ------------------------------------------------------------
\subsection{Análisis de Correlaciones y Multicolinealidad}
\label{subsec:corr}
% ------------------------------------------------------------

% ............................................................
\subsubsection{Correlaciones entre variables de engagement}
\label{sssec:corr_eng}
% ............................................................

La Figura~\ref{fig:corr_engagement} presenta la matriz de
correlaciones de Spearman entre las variables de engagement y
la dispersión entre tiempo hasta P1 y nota del Parcial~1.

\begin{figure}[htbp]
  \centering
  \includegraphics[width=\textwidth]{08_correlaciones_engagement.png}
  \caption{\textbf{Izquierda:} matriz de correlaciones de Spearman
           entre variables de engagement (triángulo inferior).
           Las celdas en rojo oscuro indican correlaciones
           altas ($r > 0{,}80$).
           \textbf{Derecha:} dispersión entre tiempo de
           engagement hasta P1 y nota del Parcial~1 por resultado
           académico, con línea de regresión y correlación de
           Spearman.}
  \label{fig:corr_engagement}
\end{figure}

Las correlaciones intra-período son altas: los pares más
problemáticos son tiempo hasta P1 $\leftrightarrow$ tiempo
hasta P2 ($r = 0{,}864$), temas~P1 $\leftrightarrow$
módulos~P1 ($r = 0{,}861$), tiempo~P2 $\leftrightarrow$
tiempo total ($r = 0{,}852$) y temas~total $\leftrightarrow$
módulos~total ($r = 0{,}941$). Esta multicolinealidad no
impide el modelado pero requiere regularización en los
modelos lineales y justifica la selección de subconjuntos
representativos de cada dimensión de engagement en lugar
de incluir todas las variables simultáneamente.

\subsubsection{Factor de inflación de varianza (VIF)}
\label{sssec:vif}

El análisis VIF sobre el conjunto de predictores candidatos
revela que las variables académicas presentan los valores más
altos (Parcial~1: VIF~$= 13{,}7$; AM~2: VIF~$= 10{,}4$;
AM~1: VIF~$= 9{,}4$), mientras que las de engagement son más
moderadas (módulos únicos P1: VIF~$= 7{,}2$;
tiempo hasta P1: VIF~$= 4{,}0$;
tiempo por visita P1: VIF~$= 3{,}9$). Ningún valor supera el
umbral de 15 que produce inestabilidad severa, pero la magnitud
de algunos VIF justifica el uso de regularización L1/L2 en los
modelos de regresión logística (ver Sección~\ref{sec:metodologia}).

% ------------------------------------------------------------
\subsection{Síntesis: Decisiones de Diseño del Dataset}
\label{subsec:sintesis_datos}
% ------------------------------------------------------------

La Tabla~\ref{tab:resumen_decisiones} sintetiza las decisiones
de preprocesamiento tomadas a partir del análisis exploratorio.

\begin{table}[htbp]
  \centering
  \caption{Resumen de decisiones de diseño del dataset con
           justificación empírica e impacto cuantificado.}
  \label{tab:resumen_decisiones}
  \small
  \begin{tabular}{p{3.5cm}p{5.0cm}p{4.8cm}}
    \toprule
    \textbf{Decisión} & \textbf{Justificación} & \textbf{Impacto} \\
    \midrule
    Excluir RETIRADO y SIN~NOTA
      & Etiqueta no observable; sesgo de retiro endógeno;
        heterogeneidad semántica
      & $N$: 764~$\to$~554 \\[3pt]
    Excluir estudiantes sin engagement
      & Cero ambiguo (75\,\% en semestre de arranque);
        $\Delta F_1 < 0{,}02$ al excluirlos
      & $N$: 554~$\to$~502 \\[3pt]
    Excluir variables de actividades evaluativas
      & AUC~$< 0{,}50$; $p > 0{,}05$ (sin discriminación)
      & $-3$ predictores candidatos \\[3pt]
    Incluir \texttt{era\_encoded}
      & $\Delta$AUC(P1) $= -0{,}087$;
        engagement mediano cae 60\,\% en era nueva
      & Control en todos los modelos \\[3pt]
    Incluir \texttt{modalidad\_am}
      & AM-3 grupal vs.\ individual:
        $p = 1{,}4\times10^{-9}$, $d = 0{,}40$
      & Control en modelos con AMs \\[3pt]
    Incluir engagement a pesar de efecto pequeño
      & $\rho_{\text{parcial}} = 0{,}166$, $p < 0{,}001$;
        valor incremental en zona marginal de P1
      & $+6$ predictores en CP1 \\[3pt]
    Aplicar balanceo SMOTE / SMOTEENN
      & 74,3\,\% APROBÓ vs.\ 25,7\,\% REPROBÓ;
        Recall sin balanceo $\approx 0{,}55$
      & Solo en training fold; nunca en evaluación \\
    \bottomrule
  \end{tabular}
\end{table}

El dataset de modelado resultante (\texttt{df\_activos\_engagement},
$N = 502$) presenta tres características estructurales que
condicionan el diseño metodológico de la siguiente sección:
(i)~\textbf{desbalance de clases} (74,3\,\% vs.\ 25,7\,\%),
que exige técnicas de remuestreo para evitar que los modelos
sesguen la predicción hacia la clase mayoritaria;
(ii)~\textbf{estructura longitudinal multisemestral}, que
invalida la validación cruzada aleatoria estándar y exige el
esquema LOSO; y (iii)~\textbf{heterogeneidad de era pedagógica},
que requiere variables de control y análisis estratificado.
Estos tres aspectos se abordan en detalle en la sección de
metodología.

% ============================================================
%  REFERENCIAS
% ============================================================

\begin{thebibliography}{9}

\bibitem{Lopez2021}
J.~López-Zambrano, J.~A. Lara~Torralbo, y C.~Romero,
``Early prediction of student learning performance through data
mining: A systematic review,''
\textit{Psicothema}, vol.~33, no.~3, pp.~456--465, 2021.

\bibitem{Tinto1993}
V.~Tinto,
\textit{Leaving College: Rethinking the Causes and Cures of
Student Attrition}, 2.ª~ed.
Chicago, IL, USA: University of Chicago Press, 1993.

\bibitem{Kerby2014}
D.~S. Kerby,
``The simple difference formula: An approach to teaching
nonparametric correlation,''
\textit{Compr. Psychol.}, vol.~3, p.~11.IT.3.4.1, 2014.

\bibitem{He2009}
H.~He y E.~A. Garcia,
``Learning from imbalanced data,''
\textit{IEEE Trans. Knowl. Data Eng.}, vol.~21, no.~9,
pp.~1263--1284, sep. 2009.

\end{thebibliography}

\end{document}
 y eliminar el
%  preámbulo y los comandos \begin{document}/\end{document}.
%  Cambiar [english] a [spanish] con texlive-lang-spanish.
%  Las figuras se cargan desde la carpeta figures/ relativa.
% ============================================================

\documentclass[12pt, letterpaper]{article}

\usepackage[utf8]{inputenc}
\usepackage[T1]{fontenc}
\usepackage[english]{babel}
\usepackage{amsmath, amssymb}
\usepackage{geometry}
\usepackage{setspace}
\usepackage{booktabs}
\usepackage{multirow}
\usepackage{array}
\usepackage{graphicx}
\usepackage{subcaption}
\usepackage{xcolor}
\usepackage{cite}
\usepackage[hidelinks]{hyperref}
\usepackage{url}
\usepackage{rotating}

\geometry{top=2.5cm, bottom=2.5cm, left=3cm, right=2.5cm}
\onehalfspacing
\renewcommand{\refname}{Referencias}

\graphicspath{{figures/}}

% ============================================================
\begin{document}
% ============================================================

\section{Descripción y Exploración de Datos}
\label{sec:exploracion}

Este capítulo caracteriza el conjunto de datos empleado en el SAT,
documenta las decisiones de preprocesamiento adoptadas y presenta
los hallazgos del análisis exploratorio que fundamentan el diseño
del sistema de modelado. La presentación se organiza en seis
bloques: (i)~fuente y composición del dataset;
(ii)~definición y justificación de los subconjuntos de trabajo;
(iii)~calidad y estructura de los datos;
(iv)~análisis estadístico de las notas académicas;
(v)~análisis de la actividad en plataforma; y
(vi)~análisis de correlaciones y síntesis de decisiones de diseño.

% ------------------------------------------------------------
\subsection{Fuente y Composición del Conjunto de Datos}
\label{subsec:fuente}
% ------------------------------------------------------------

Los datos provienen del curso IIND-2104 \textit{Modelos
Probabilísticos} de la Universidad de los Andes, dictado entre
los semestres 202320 y 202520 —cinco cohortes consecutivas
cubiertas entre 2023 y 2025. La información se integra a partir
de dos fuentes complementarias:

\begin{itemize}
  \item \textbf{Sistema de gestión académica institucional:}
    calificaciones oficiales de los tres parciales y de las
    evaluaciones continuas (actividades modulares, quizzes,
    talleres y proyectos, según el formato de cada semestre).

  \item \textbf{Plataforma Bloque Neón:} sistema de gestión del
    aprendizaje (LMS) que registra, por sesión y por módulo
    de contenido, el tiempo de estudio, los temas consultados,
    los módulos visitados y la frecuencia de acceso de cada
    estudiante.
\end{itemize}

El dataset crudo consolidado cuenta con \textbf{764~registros y
77~variables}, con cuatro posibles resultados académicos finales:
\texttt{APROBÓ} (413 estudiantes, 54,1\,\%),
\texttt{REPROBÓ} (141, 18,5\,\%),
\texttt{RETIRADO} (164, 21,5\,\%) y
\texttt{SIN~NOTA} (46, 6,0\,\%).

La Figura~\ref{fig:composicion} ilustra cómo se distribuyen
estos resultados por semestre en términos absolutos y
porcentuales, y la Tabla~\ref{tab:composicion_semestral} detalla
las frecuencias exactas.

\begin{figure}[htbp]
  \centering
  \includegraphics[width=\textwidth]{01_composicion_dataset.png}
  \caption{Composición del dataset por semestre y resultado
           académico. \textbf{Izquierda:} distribución absoluta.
           \textbf{Derecha:} distribución porcentual con
           etiquetas de proporción. El sombreado naranja
           corresponde a la era pedagógica antigua (202320--202410)
           y el verde a la nueva (202420--202520).}
  \label{fig:composicion}
\end{figure}

\begin{table}[htbp]
  \centering
  \caption{Composición por semestre, era pedagógica y resultado
           académico final.}
  \label{tab:composicion_semestral}
  \small
  \begin{tabular}{lccrrrrrrr}
    \toprule
    \textbf{Semestre} & \textbf{Era} & \textbf{Total}
      & \multicolumn{2}{c}{\textbf{APROBÓ}}
      & \multicolumn{2}{c}{\textbf{REPROBÓ}}
      & \multicolumn{2}{c}{\textbf{RETIRADO}}
      & \textbf{SIN} \\
    \cmidrule(lr){4-5}\cmidrule(lr){6-7}\cmidrule(lr){8-9}
    & & & $n$ & \% & $n$ & \% & $n$ & \% & \textbf{NOTA} \\
    \midrule
    202320 & Antigua & 170 & 116 & 68,2 &  27 & 15,9 &  27 & 15,9 & 0  \\
    202410 & Antigua & 150 &  68 & 45,3 &  14 &  9,3 &  30 & 20,0 & 38 \\
    202420 & Nueva   & 164 &  90 & 54,9 &  33 & 20,1 &  40 & 24,4 & 1  \\
    202510 & Nueva   & 140 &  70 & 50,0 &  47 & 33,6 &  21 & 15,0 & 2  \\
    202520 & Nueva   & 140 &  69 & 49,3 &  20 & 14,3 &  46 & 32,9 & 5  \\
    \midrule
    \textbf{Total} & — & \textbf{764}
      & \textbf{413} & \textbf{54,1}
      & \textbf{141} & \textbf{18,5}
      & \textbf{164} & \textbf{21,5}
      & \textbf{46} \\
    \bottomrule
  \end{tabular}
\end{table}

Dos fenómenos merecen atención especial. Primero, la tasa de
reprobación en 202510 (33,6\,\%) es la más alta del período y
más del doble de la de 202410 (9,3\,\%), lo que sugiere
variabilidad intercohorte que el modelo debe generalizar. Segundo,
el semestre 202520 registra la mayor tasa de retiro (32,9\,\%),
fenómeno atípico que podría obedecer a cambios en la política
institucional de permanencia o a condiciones externas no
capturadas en los datos.

% ............................................................
\subsubsection{Dos eras pedagógicas}
\label{sssec:eras}
% ............................................................

Un rasgo estructural del dataset es la coexistencia de dos
\textbf{formatos de evaluación continua}:

\begin{itemize}
  \item \textbf{Era Antigua} (202320 y 202410; $n = 320$ estudiantes
    activos): las evaluaciones continuas consisten en talleres
    individuales y proyectos por fases grupales.

  \item \textbf{Era Nueva} (202420, 202510 y 202520; $n = 444$
    estudiantes activos): las evaluaciones continuas son
    actividades modulares (AMs) y quizzes individuales.
\end{itemize}

Dentro de la era nueva existe una distinción adicional: en
202420 las AMs se realizaron \textbf{en grupo}, mientras que
a partir de 202510 pasaron a ser \textbf{individuales}. Este
cambio produce diferencias estadísticamente significativas en la
distribución de AM-3 ($U$-test, $p = 1{,}39 \times 10^{-9}$,
$d_{\text{Cohen}} = 0{,}40$, AUC$_{\text{grupo}} = 0{,}369$
vs.\ AUC$_{\text{ind.}} = 0{,}246$), lo que justifica incluir la
variable \texttt{modalidad\_am} como control en los modelos
entrenados con datos del formato nuevo.

% ------------------------------------------------------------
\subsection{Definición y Justificación de los Subconjuntos
            de Trabajo}
\label{subsec:subsets}
% ------------------------------------------------------------

El dataset bruto requiere dos depuraciones sucesivas antes de
ser apto para el modelado de clasificación.

% ............................................................
\subsubsection{Exclusión de estudiantes retirados y sin
               calificación}
\label{sssec:excl_retirados}
% ............................................................

Los 164 estudiantes con resultado \texttt{RETIRADO} y los
46 con \texttt{SIN~NOTA} se excluyen por tres razones:

\begin{enumerate}
  \item \textbf{Ausencia de etiqueta objetivo observable.}
    El SAT modela $P(\text{REPROBÓ} \mid \mathbf{x}_i)$, la
    probabilidad de que un estudiante \emph{finalice} el semestre
    reprobando. Los estudiantes retirados abandonaron el curso
    antes de producirse el resultado; su etiqueta final es
    desconocida y no puede inferirse sin supuestos fuertes.

  \item \textbf{Sesgo de retiro endógeno.}
    La decisión de retirarse está correlacionada con el
    desempeño anticipado: los estudiantes que perciben que van
    a reprobar se retiran en mayor proporción \cite{Tinto1993}.
    Incluirlos en la clase \texttt{REPROBÓ} introduciría sesgo
    de selección en las variables de engagement: un estudiante
    retirado en la semana~4 tiene engagement mínimo no por falta
    de habilidad sino por abandono temprano, contaminando el
    predictor más utilizado.

  \item \textbf{Heterogeneidad semántica.}
    La predicción del retiro como fenómeno independiente
    requeriría modelos multi-clase y supuestos adicionales
    sobre su equivalencia con la reprobación, lo cual está
    fuera del alcance del presente trabajo.
\end{enumerate}

El subconjunto resultante es \textbf{\texttt{df\_activos}}
($N = 554$; APROBÓ: 413, REPROBÓ: 141).

% ............................................................
\subsubsection{Exclusión adicional de estudiantes sin actividad
               en plataforma}
\label{sssec:excl_engagement}
% ............................................................

Del subconjunto \texttt{df\_activos}, 52 estudiantes (9,4\,\%)
presentan \textbf{engagement total nulo} en Bloque Neón durante
todo el semestre (i.e., \texttt{engagement\_hasta\_p1} $= 0$).
Estos casos se excluyen del subconjunto de modelado por las
razones siguientes:

\begin{enumerate}
  \item \textbf{Ambigüedad del cero.}
    Existen dos interpretaciones mutuamente excluyentes:
    (a)~el estudiante genuinamente no utilizó la plataforma,
    o (b)~ocurrió un error en el pipeline de integración de
    datos. Ambas hipótesis son plausibles: 39 de los 52 casos
    (75\,\%) se concentran en el semestre 202320, el primero
    en incorporar Bloque Neón, cuando el sistema estaba aún
    en fase de implementación. No es posible distinguir entre
    causas con la información disponible.

  \item \textbf{Riesgo de imputación espuria.}
    Imputar cero como valor real de engagement asumiría que
    ``no usar la plataforma'' es una señal homogénea entre
    semestres, cuando los patrones de no-uso varían
    considerablemente entre cohortes.

  \item \textbf{Impacto estadístico marginal confirmado.}
    Se realizó un análisis de sensibilidad: los modelos
    entrenados sobre \texttt{df\_activos} completo ($N = 554$)
    y sobre el subconjunto filtrado ($N = 502$) arrojan
    diferencias en F1 menores a $\Delta F_1 < 0{,}02$,
    confirmando que la exclusión no altera las conclusiones.
\end{enumerate}

El subconjunto de modelado final es
\textbf{\texttt{df\_activos\_engagement}}
($N = 502$; APROBÓ:~373, REPROBÓ:~129).
La Tabla~\ref{tab:subsets} resume los tres subconjuntos
empleados en el análisis.

\begin{table}[htbp]
  \centering
  \caption{Subconjuntos de trabajo. \texttt{df\_activos\_engagement}
           es el conjunto empleado para el entrenamiento y
           evaluación de todos los modelos del SAT.}
  \label{tab:subsets}
  \begin{tabular}{lrrrl}
    \toprule
    \textbf{Subconjunto} & \textbf{$N$}
      & \textbf{APROBÓ} & \textbf{REPROBÓ}
      & \textbf{Uso principal} \\
    \midrule
    \texttt{df\_all}
      & 764 & 413 & 141
      & Tasas globales, análisis de retiro \\
    \texttt{df\_activos}
      & 554 & 413 & 141
      & Análisis de notas y distribuciones \\
    \texttt{df\_activos\_engagement}
      & \textbf{502} & \textbf{373} & \textbf{129}
      & \textbf{Entrenamiento y evaluación de modelos} \\
    \bottomrule
  \end{tabular}
\end{table}

% ------------------------------------------------------------
\subsection{Calidad y Estructura de los Datos}
\label{subsec:calidad}
% ------------------------------------------------------------

% ............................................................
\subsubsection{Nulos estructurales versus problemáticos}
\label{sssec:nulos}
% ............................................................

El dataset presenta un volumen considerable de valores faltantes,
en su mayoría atribuibles al cambio de formato pedagógico.
Se distinguen dos categorías.

\paragraph{Nulos estructurales (esperados).}
Son variables que solo existen en una era y, por tanto, están
ausentes de forma legítima en la otra:
\texttt{quiz\_4/5/6} (77,8\,\% de nulos, solo en algunos
semestres del formato nuevo),
\texttt{taller\_1/2/3} (59,4\,\%, exclusivo de la era antigua),
\texttt{fase\_1/2} (59,4\,\%, proyectos grupales de la era antigua),
y \texttt{am\_4/5/6} (56,7\,\%, AMs opcionales no presentes en
todos los semestres). Estos nulos no representan problemas de
calidad sino de diseño del instrumento de evaluación.

\paragraph{Nulos problemáticos (casos especiales).}
Son faltantes en variables que deberían estar presentes para
todos los estudiantes activos:
\begin{itemize}
  \item \textbf{Parcial~1 $= 0$ ($n = 1$):} un único estudiante
    registra nota cero. No es posible distinguir ausencia de nota
    real. Se retiene con variable indicadora
    \texttt{no\_presento\_p1}.
  \item \textbf{Parcial~3 $= 0$ ($n = 15$; 13 reprobaron):}
    concentrado en estudiantes encaminados a reprobar, consistente
    con ausentismo por desmotivación al final del semestre.
    Se retienen con su valor real.
  \item \textbf{Notas superiores a 5{,}0:} presentes por
    bonificaciones otorgadas por docentes. Son valores legítimos
    y se retienen sin transformación.
\end{itemize}

% ............................................................
\subsubsection{Variables de plataforma sin poder discriminativo}
\label{sssec:excl_vars}
% ............................................................

La plataforma Bloque Neón registra tanto métricas de
\textit{estudio de contenidos} como métricas de
\textit{actividades evaluativas en línea}. El análisis estadístico
univariante sobre estas últimas —\texttt{actividades\_n\_realizadas},
\texttt{actividades\_tiempo\_total} y
\texttt{actividades\_visitas\_total}— revela ausencia de
poder discriminativo entre aprobados y reprobados
(Tabla~\ref{tab:excl_vars}), lo que se ilustra también en la
Figura~\ref{fig:vars_excluidas}.

\begin{table}[htbp]
  \centering
  \caption{Variables de actividades evaluativas de la plataforma.
           AUC~$\approx 0{,}5$ y $p > 0{,}05$ indican que no
           discriminan mejor que la clasificación aleatoria.}
  \label{tab:excl_vars}
  \begin{tabular}{lrrrr}
    \toprule
    \textbf{Variable}
      & \textbf{Med.\ APROBÓ}
      & \textbf{Med.\ REPROBÓ}
      & \textbf{AUC}
      & \textbf{$p$-valor (MW)} \\
    \midrule
    \texttt{actividades\_n\_realizadas}   & 1,00 & 1,00 & 0,473 & 0,830 \\
    \texttt{actividades\_tiempo\_total}   & 1,15 & 1,10 & 0,498 & 0,532 \\
    \texttt{actividades\_visitas\_total}  & 9,00 & 11,00 & 0,465 & 0,886 \\
    \bottomrule
  \end{tabular}
\end{table}

\begin{figure}[htbp]
  \centering
  \includegraphics[width=\textwidth]{11_variables_excluidas.png}
  \caption{Distribuciones de las tres variables de actividades
           evaluativas para aprobados (azul) y reprobados (rojo).
           Los valores de $p > 0{,}05$ y AUC~$< 0{,}5$
           confirman la ausencia de poder discriminativo,
           justificando su exclusión del conjunto de predictores.}
  \label{fig:vars_excluidas}
\end{figure}

Nótese que, para \texttt{actividades\_visitas\_total}, la mediana
de los reprobados (11) supera la de los aprobados (9), con
AUC~$= 0{,}465$ —peor que la clasificación aleatoria—, lo que
sugiere que los estudiantes con mayor número de visitas a
actividades no necesariamente obtienen mejores resultados. Las
tres variables son excluidas del conjunto de predictores para
evitar introducir ruido en los modelos.

% ------------------------------------------------------------
\subsection{Análisis Estadístico de las Notas Académicas}
\label{subsec:notas}
% ------------------------------------------------------------

El análisis de notas se realiza sobre \texttt{df\_activos}
($N = 554$) para preservar la distribución real sin los
filtros de engagement.

% ............................................................
\subsubsection{Distribución de la nota final}
\label{sssec:dist_nota}
% ............................................................

La nota final (escala 0--5) presenta una distribución bimodal
con dos concentraciones claramente separadas: una en torno a
2,6 (zona de reprobación) y otra en torno a 3,7 (zona de
aprobación cómoda), como se observa en la
Figura~\ref{fig:dist_nota}.

\begin{figure}[htbp]
  \centering
  \includegraphics[width=\textwidth]{02_distribucion_nota_final.png}
  \caption{Distribución de la nota final.
           \textbf{Izquierda:} histogramas superpuestos con
           curva de densidad KDE por resultado académico.
           La línea discontinua vertical indica el umbral de
           aprobación (3{,}0).
           \textbf{Derecha:} diagramas de caja por semestre y
           resultado; el sombreado distingue las eras pedagógicas.}
  \label{fig:dist_nota}
\end{figure}

Las estadísticas descriptivas detalladas se presentan en la
Tabla~\ref{tab:desc_nota}. La prueba de Mann-Whitney confirma
la separación altamente significativa entre grupos
($U = 58\,233{,}0$, $p = 9{,}97 \times 10^{-71}$,
$d_{\text{Cohen}} = 2{,}41$). Esta magnitud del efecto implica
que la distribución de notas de los aprobados y la de los
reprobados se solapan apenas en una pequeña región central
(alrededor de $[2{,}8;\;3{,}2]$), que es precisamente la zona
donde los modelos de clasificación concentran su incertidumbre.

\begin{table}[htbp]
  \centering
  \caption{Estadísticas descriptivas de la nota final por resultado
           y era pedagógica. $\bar{x}$: media; $\tilde{x}$: mediana;
           $s$: desviación estándar.}
  \label{tab:desc_nota}
  \begin{tabular}{lrrrr}
    \toprule
    \textbf{Grupo} & \textbf{$N$}
      & \textbf{$\bar{x}$} & \textbf{$\tilde{x}$} & \textbf{$s$} \\
    \midrule
    APROBÓ (global)                & 413 & 3,724 & 3,620 & 0,527 \\
    REPROBÓ (global)               & 141 & 2,418 & 2,632 & 0,583 \\
    \midrule
    Era Antigua — APROBÓ           & 184 & 3,777 & 3,680 & 0,504 \\
    Era Antigua — REPROBÓ          &  41 & 2,446 & 2,720 & 0,671 \\
    Era Nueva   — APROBÓ           & 229 & 3,682 & 3,530 & 0,541 \\
    Era Nueva   — REPROBÓ          & 100 & 2,407 & 2,605 & 0,545 \\
    \bottomrule
  \end{tabular}
\end{table}

% ............................................................
\subsubsection{Poder predictivo de los parciales:
               análisis por era pedagógica}
\label{sssec:parciales_era}
% ............................................................

La Figura~\ref{fig:parciales_era} compara la distribución de
las tres calificaciones de parciales entre aprobados y reprobados,
estratificada por era pedagógica.

\begin{figure}[htbp]
  \centering
  \includegraphics[width=\textwidth]{05_parciales_era_resultado.png}
  \caption{Distribuciones de Parcial~1, Parcial~2 y Parcial~3
           por resultado académico, estratificadas por era
           pedagógica. Cada subgráfico incluye el $p$-valor
           de la prueba Mann-Whitney, el AUC univariante y el
           coeficiente de rango biserial~$r_b$.}
  \label{fig:parciales_era}
\end{figure}

La Tabla~\ref{tab:auc_era} sintetiza los AUC univariantes por
variable y era. El hallazgo más relevante es la \textbf{reducción
del poder predictivo de P1 en la era nueva}
($\Delta\text{AUC} = -0{,}087$). Esta caída puede atribuirse
al mayor peso relativo de las actividades continuas (quizzes y
AMs) en el nuevo currículo, que diluye la señal del primer
examen. Contrariamente, el Parcial~2 mantiene una capacidad
predictiva alta y estable en ambas eras
($\Delta\text{AUC} = -0{,}025$), lo que refuerza el valor
del segundo checkpoint (CP2) del SAT.

\begin{table}[htbp]
  \centering
  \caption{AUC univariante de las variables académicas por era
           pedagógica. $\Delta$AUC indica la variación entre eras.
           Diferencias $> 5\,\%$ se consideran sustanciales.}
  \label{tab:auc_era}
  \begin{tabular}{lrrrc}
    \toprule
    \textbf{Variable}
      & \textbf{AUC Era Antigua}
      & \textbf{AUC Era Nueva}
      & \textbf{$\Delta$AUC}
      & \textbf{Sustancial} \\
    \midrule
    Parcial~1      & 0,869 & 0,781 & $-0{,}087$ & Sí \\
    Parcial~2      & 0,925 & 0,900 & $-0{,}025$ & No \\
    Parcial~3      & 0,931 & 0,845 & $-0{,}086$ & Sí \\
    Prom.\ AMs     & 0,737 & 0,731 & $-0{,}006$ & No \\
    \bottomrule
  \end{tabular}
\end{table}

La Figura~\ref{fig:eras_comparacion} ilustra visualmente esta
comparación, tanto para las variables académicas como para los
patrones de engagement entre eras.

\begin{figure}[htbp]
  \centering
  \includegraphics[width=\textwidth]{09_comparacion_eras.png}
  \caption{\textbf{Izquierda:} AUC univariante de variables
           académicas por era pedagógica. La línea discontinua
           roja marca el umbral mínimo del SAT (AUC~$= 0{,}75$).
           \textbf{Derecha:} medianas de engagement por era
           y resultado académico, normalizadas para comparación
           entre variables de diferente escala.}
  \label{fig:eras_comparacion}
\end{figure}

% ............................................................
\subsubsection{Poder predictivo temprano del Parcial~1:
               tabla de riesgo empírico}
\label{sssec:riesgo_p1}
% ............................................................

El primer checkpoint del SAT (CP1, semana~6) opera únicamente
con la información disponible en ese momento, siendo el Parcial~1
el predictor individual más potente. La
Figura~\ref{fig:riesgo_p1} y la Tabla~\ref{tab:riesgo_p1}
cuantifican el riesgo empírico de reprobación por intervalo
de nota de P1 sobre \texttt{df\_activos} ($N = 554$).

\begin{figure}[htbp]
  \centering
  \includegraphics[width=\textwidth]{03_riesgo_p1.png}
  \caption{\textbf{Izquierda:} diagrama de dispersión entre la
           nota del Parcial~1 y la nota final, coloreado por
           resultado académico. Las líneas discontinuas marcan
           el umbral de aprobación (3{,}0). La correlación de
           Spearman se muestra en el recuadro.
           \textbf{Derecha:} tasa de reprobación empírica por
           intervalo de P1. La línea discontinua horizontal
           indica la tasa global de reprobación (25,7\,\%).}
  \label{fig:riesgo_p1}
\end{figure}

\begin{table}[htbp]
  \centering
  \caption{Tabla de riesgo empírico de reprobación por intervalo
           de nota en Parcial~1 (CP1, semana~6).}
  \label{tab:riesgo_p1}
  \begin{tabular}{ccrrrr}
    \toprule
    \textbf{Intervalo P1}
      & \textbf{Nivel}
      & \textbf{$N$}
      & \textbf{P(APROBÓ)}
      & \textbf{P(REPROBÓ)}
      & \textbf{$\bar{x}$ nota final} \\
    \midrule
    $[0{,}0;\;1{,}5)$  & Crítico  &  19 & 15,8\,\% & \textbf{84,2\,\%} & 2,16 \\
    $[1{,}5;\;2{,}5)$  & Alto     & 102 & 43,1\,\% & \textbf{56,9\,\%} & 2,80 \\
    $[2{,}5;\;3{,}0)$  & Moderado & 100 & 68,0\,\% &         32,0\,\%  & 3,10 \\
    $[3{,}0;\;3{,}5)$  & Límite   & 128 & 85,9\,\% &         14,1\,\%  & 3,44 \\
    $[3{,}5;\;4{,}0)$  & Seguro   &  95 & 83,2\,\% &         16,8\,\%  & 3,62 \\
    $[4{,}0;\;5{,}0]$  & Sólido   & 110 & 99,1\,\% &          0,9\,\%  & 4,18 \\
    \bottomrule
  \end{tabular}
\end{table}

Los resultados son contundentes en los extremos: estudiantes con
P1 $< 1{,}5$ tienen una probabilidad empírica del 84,2\,\% de
reprobar, mientras que con P1~$\geq 4{,}0$ esa probabilidad
cae al 0,9\,\%. Sin embargo, el intervalo $[2{,}5;\;3{,}5)$,
que concentra 228 estudiantes (41,2\,\% del total activo),
presenta tasas de reprobación intermedias (14--32\,\%) donde la
nota del parcial tiene menor poder de discriminación. Este es
precisamente el segmento donde el engagement y otros predictores
complementarios tienen mayor valor incremental.

% ............................................................
\subsubsection{Trayectorias académicas P1~$\to$~P2~$\to$~P3}
\label{sssec:trayectorias}
% ............................................................

Más allá del nivel absoluto de P1, la \emph{trayectoria}
de desempeño a lo largo de los tres parciales provee
información adicional sobre el resultado final. Se identifican
cinco patrones cualitativamente distintos
(Figura~\ref{fig:trayectorias} y Tabla~\ref{tab:trayectorias}).

\begin{figure}[htbp]
  \centering
  \includegraphics[width=\textwidth]{04_trayectorias_parciales.png}
  \caption{Trayectorias de desempeño académico entre parciales.
           \textbf{Izquierda:} frecuencia de cada trayectoria
           por resultado académico.
           \textbf{Derecha:} evolución de la nota promedio
           en cada parcial por trayectoria, con el porcentaje
           de aprobación de cada grupo en la leyenda.
           La línea discontinua marca el umbral de 3{,}0.}
  \label{fig:trayectorias}
\end{figure}

\begin{table}[htbp]
  \centering
  \caption{Patrones de trayectoria académica identificados en
           la secuencia P1~$\to$~P2~$\to$~P3. Los porcentajes
           se calculan dentro de cada grupo de trayectoria.}
  \label{tab:trayectorias}
  \begin{tabular}{lrrrr}
    \toprule
    \textbf{Trayectoria} & \textbf{Total}
      & \textbf{APROBÓ} & \textbf{REPROBÓ} & \textbf{\% APR.} \\
    \midrule
    Declive sostenido (P1~$>$~P2~$>$~P3)    &  40 &  22 &  18 & 55,0\,\% \\
    Estable   (P1~$\approx$~P2~$\approx$~P3)& 306 & 225 &  81 & 73,5\,\% \\
    Mejora sostenida (P1~$<$~P2~$<$~P3)     &  22 &  20 &   2 & 90,9\,\% \\
    Pico intermedio (P1~$<$~P2~$>$~P3)      & 129 &  99 &  30 & 76,7\,\% \\
    Recuperación post-P1 (P1~$<$~P2)        &  57 &  47 &  10 & 82,5\,\% \\
    \bottomrule
  \end{tabular}
\end{table}

El grupo de \textbf{declive sostenido} concentra la mayor tasa
de reprobación (45,0\,\%) y es identificable a partir de la
comparación P1~$\to$~P2 disponible en CP2. La
\textbf{mejora sostenida} es el patrón con mayor probabilidad
de aprobación (90,9\,\%), lo que indica que la trayectoria
ascendente es una señal de éxito más fiable que el nivel
absoluto de P1.

Más relevante aún para el diseño del SAT: 115 de los 221
estudiantes con P1~$< 3{,}0$ terminaron aprobando (52,0\,\%),
con nota media de Parcial~2 $= 3{,}46$ y nota final media
$= 3{,}38$. Este fenómeno de \textit{recuperación académica}
demuestra que un clasificador basado únicamente en P1 tendría
una tasa de falsas alarmas del 52\,\% entre los estudiantes de
riesgo moderado, motivando explícitamente la revisión del estado
de alerta en el segundo checkpoint (CP2, semana~11).

% ------------------------------------------------------------
\subsection{Análisis de la Actividad en Plataforma
            (Engagement)}
\label{subsec:eng}
% ------------------------------------------------------------

El análisis de engagement se realiza sobre
\texttt{df\_activos\_engagement} ($N = 502$): 373 aprobados y
129 reprobados. La Figura~\ref{fig:composicion_modelado}
muestra la composición de este subconjunto.

\begin{figure}[htbp]
  \centering
  \includegraphics[width=0.92\textwidth]{12_composicion_modelado.png}
  \caption{Composición del subconjunto de modelado
           (\texttt{df\_activos\_engagement}, $N = 502$).
           \textbf{Izquierda:} distribución de clases. El
           desbalance 74,3\,\%~/ 25,7\,\% motiva el uso de
           técnicas de balanceo (SMOTE, SMOTEENN).
           \textbf{Derecha:} descomposición por semestre, con
           el porcentaje de reprobados sobre cada barra.}
  \label{fig:composicion_modelado}
\end{figure}

% ............................................................
\subsubsection{Descripción de las métricas de engagement}
\label{sssec:vars_eng}
% ............................................................

Bloque Neón registra el comportamiento de estudio en dos
dimensiones complementarias:

\begin{itemize}
  \item \textbf{Amplitud de estudio (cobertura):}
    módulos únicos visitados (\texttt{parcial\_X\_modulos\_unicos}),
    temas únicos consultados (\texttt{parcial\_X\_temas\_unicos})
    y visitas totales al contenido (\texttt{parcial\_X\_visitas}).

  \item \textbf{Profundidad de estudio (intensidad):}
    tiempo promedio por visita
    (\texttt{parcial\_X\_tiempo\_por\_visita}) y visitas promedio
    por tema (\texttt{parcial\_X\_visitas\_por\_tema}).
\end{itemize}

Además, se dispone del tiempo acumulado hasta cada checkpoint
(\texttt{engagement\_hasta\_p1},
\texttt{engagement\_hasta\_p2}) y del total del semestre
(\texttt{total\_tiempo\_hrs}).

% ............................................................
\subsubsection{Diferencias estadísticas entre aprobados y
               reprobados: prueba de Mann-Whitney}
\label{sssec:mann_eng}
% ............................................................

La Figura~\ref{fig:engagement_boxplots} presenta las
distribuciones de las principales variables de engagement
disponibles en CP1 por semestre y resultado, y la
Tabla~\ref{tab:mw_engagement} reporta los estadísticos
completos de las pruebas Mann-Whitney.

\begin{figure}[htbp]
  \centering
  \includegraphics[width=\textwidth]{07_engagement_boxplots.png}
  \caption{Distribución de las cuatro principales variables de
           engagement en CP1 por semestre y resultado académico.
           El sombreado naranja indica la era antigua y el verde
           la nueva. En cada subgráfico se reportan el $p$-valor
           de Mann-Whitney, el AUC univariante y el coeficiente
           de rango biserial $r_b$.}
  \label{fig:engagement_boxplots}
\end{figure}

\begin{table}[htbp]
  \centering
  \caption{Prueba Mann-Whitney para variables de engagement
           disponibles en CP1. $\tilde{x}$: mediana; $r_b$:
           coeficiente de rango biserial \cite{Kerby2014}.
           Significancia: *** $p < 0{,}001$; * $p < 0{,}05$.}
  \label{tab:mw_engagement}
  \begin{tabular}{lrrrrc}
    \toprule
    \textbf{Variable}
      & \textbf{$\tilde{x}$ APR.}
      & \textbf{$\tilde{x}$ REP.}
      & \textbf{AUC}
      & \textbf{$r_b$}
      & \textbf{$p$-valor} \\
    \midrule
    Tiempo hasta P1 (h)
      & 8,27  & 5,71  & 0,602 & 0,204
      & $2{,}8\times10^{-4}$ *** \\
    Módulos únicos P1
      & 7,00  & 6,00  & 0,601 & 0,203
      & $2{,}7\times10^{-4}$ *** \\
    Visitas contenido P1
      & 37,00 & 32,00 & 0,575 & 0,150
      & $1{,}1\times10^{-2}$ *   \\
    Temas únicos P1
      & 12,00 & 11,00 & 0,567 & 0,134
      & $2{,}3\times10^{-2}$ *   \\
    Visitas por tema P1
      & 3,20  & 2,83  & 0,565 & 0,130
      & $2{,}8\times10^{-2}$ *   \\
    Minutos por visita P1
      & 12,26 & 11,07 & 0,562 & 0,125
      & $3{,}5\times10^{-2}$ *   \\
    \bottomrule
  \end{tabular}
\end{table}

Todas las variables de engagement en CP1 presentan diferencias
estadísticamente significativas ($p < 0{,}05$), pero los
tamaños del efecto son \textbf{pequeños}
($r_b \in [0{,}12;\;0{,}20]$), muy por debajo de los observados
para las variables académicas (Tabla~\ref{tab:mw_academicas}).

\begin{table}[htbp]
  \centering
  \caption{Prueba Mann-Whitney para variables académicas.
           Todas presentan tamaños del efecto grandes
           ($d_{\text{Cohen}} \geq 0{,}91$)
           con significancia $p < 10^{-15}$.}
  \label{tab:mw_academicas}
  \begin{tabular}{lrrrrrr}
    \toprule
    \textbf{Variable}
      & \textbf{$\tilde{x}$ APR.}
      & \textbf{$\tilde{x}$ REP.}
      & \textbf{AUC}
      & \textbf{$r_b$}
      & \textbf{$d_{\text{Cohen}}$}
      & \textbf{$p$-valor} \\
    \midrule
    Parcial~1       & 3,42 & 2,43 & 0,810 & 0,620 & 1,24
      & $1{,}8\times10^{-28}$ \\
    Parcial~2       & 3,80 & 2,33 & 0,861 & 0,722 & 1,56
      & $6{,}7\times10^{-38}$ \\
    Parcial~3       & 3,48 & 1,70 & 0,882 & 0,763 & 1,66
      & $4{,}7\times10^{-42}$ \\
    Prom.\ AMs      & 4,11 & 3,12 & 0,728 & 0,456 & 0,91
      & $2{,}9\times10^{-16}$ \\
    Prom.\ Quizzes  & 4,00 & 2,91 & 0,821 & 0,642 & 1,36
      & $1{,}0\times10^{-20}$ \\
    \bottomrule
  \end{tabular}
\end{table}

% ............................................................
\subsubsection{Poder discriminativo individual: ranking de AUC}
\label{sssec:auc_ranking}
% ............................................................

La Figura~\ref{fig:auc_univariante} resume el poder discriminativo
individual de cada variable disponible en CP1 mediante su AUC
univariante, ordenado de mayor a menor.

\begin{figure}[htbp]
  \centering
  \includegraphics[width=0.92\textwidth]{06_auc_univariante_cp1.png}
  \caption{AUC univariante de cada predictor candidato para CP1.
           Las variables académicas (azul) dominan ampliamente;
           las de engagement (verde) se ubican en la zona
           $[0{,}56;\;0{,}60]$, superando la clasificación
           aleatoria pero por debajo del umbral mínimo del SAT
           (0{,}75, línea verde discontinua). Prom.\ Quizzes no
           está disponible en CP1 para todos los semestres.}
  \label{fig:auc_univariante}
\end{figure}

La brecha entre el AUC de Parcial~1 (0,810) y el del mejor
predictor de engagement (Tiempo hasta P1: 0,602) es de 20,8
puntos porcentuales. Esto confirma que el engagement es un
predictor \emph{complementario}, no un sustituto de las
calificaciones. Su valor radica en los casos marginales donde
la nota de P1 tiene menor poder de discriminación
(zona $[2{,}5;\;3{,}5)$ de la tabla de riesgo).

% ............................................................
\subsubsection{Valor incremental del engagement y correlación
               parcial}
\label{sssec:eng_incremental}
% ............................................................

Para cuantificar el aporte del engagement \emph{más allá} de lo
que ya captura Parcial~1, se calcula la correlación de Spearman
entre el tiempo de engagement hasta P1 y la nota final,
controlando por la nota de P1. El resultado es
$\rho_{\text{parcial}} = 0{,}166$ ($p = 1{,}9 \times 10^{-4}$,
$N = 502$), indicando que el engagement explica un 2--3\,\% de
varianza adicional en el resultado. Aunque modesto en términos
absolutos, este aporte es estadísticamente robusto y
operacionalmente relevante en la zona de incertidumbre del
clasificador.

% ............................................................
\subsubsection{Patrones de acumulación de engagement entre
               checkpoints}
\label{sssec:eng_evolucion}
% ............................................................

La Figura~\ref{fig:engagement_evolucion} explora cómo se acumula
el engagement entre los dos checkpoints del SAT.

\begin{figure}[htbp]
  \centering
  \includegraphics[width=\textwidth]{10_evolucion_engagement.png}
  \caption{\textbf{Izquierda:} correlación entre el tiempo de
           engagement hasta P1 y hasta P2 por resultado académico
           ($\rho_{\text{Spearman}}$ reportado). La diagonal
           discontinua indica el punto de no-incremento entre
           checkpoints.
           \textbf{Derecha:} distribución del incremento de
           engagement en cada período por resultado (violines).
           Los $p$-valores corresponden a pruebas Mann-Whitney
           entre estados.}
  \label{fig:engagement_evolucion}
\end{figure}

La correlación entre el engagement hasta P1 y hasta P2 es alta
para ambos grupos (aprobados: $\rho \approx 0{,}86$;
reprobados: similar), lo que indica \textbf{persistencia de los
hábitos de estudio}: los estudiantes que estudian más en el primer
período tienden a hacerlo también en el segundo. Los aprobados
acumulan en el período P1~$\to$~P2 una mediana de 6,80 horas
adicionales frente a 4,85 de los reprobados; esta diferencia
se amplía hacia el final del semestre (5,29~h vs.~3,28~h en
el período P2~$\to$~fin). Este patrón sugiere que los estudiantes
en riesgo no solo estudian menos, sino que además reducen
progresivamente su dedicación conforme avanza el semestre.

% ------------------------------------------------------------
\subsection{Análisis de Correlaciones y Multicolinealidad}
\label{subsec:corr}
% ------------------------------------------------------------

% ............................................................
\subsubsection{Correlaciones entre variables de engagement}
\label{sssec:corr_eng}
% ............................................................

La Figura~\ref{fig:corr_engagement} presenta la matriz de
correlaciones de Spearman entre las variables de engagement y
la dispersión entre tiempo hasta P1 y nota del Parcial~1.

\begin{figure}[htbp]
  \centering
  \includegraphics[width=\textwidth]{08_correlaciones_engagement.png}
  \caption{\textbf{Izquierda:} matriz de correlaciones de Spearman
           entre variables de engagement (triángulo inferior).
           Las celdas en rojo oscuro indican correlaciones
           altas ($r > 0{,}80$).
           \textbf{Derecha:} dispersión entre tiempo de
           engagement hasta P1 y nota del Parcial~1 por resultado
           académico, con línea de regresión y correlación de
           Spearman.}
  \label{fig:corr_engagement}
\end{figure}

Las correlaciones intra-período son altas: los pares más
problemáticos son tiempo hasta P1 $\leftrightarrow$ tiempo
hasta P2 ($r = 0{,}864$), temas~P1 $\leftrightarrow$
módulos~P1 ($r = 0{,}861$), tiempo~P2 $\leftrightarrow$
tiempo total ($r = 0{,}852$) y temas~total $\leftrightarrow$
módulos~total ($r = 0{,}941$). Esta multicolinealidad no
impide el modelado pero requiere regularización en los
modelos lineales y justifica la selección de subconjuntos
representativos de cada dimensión de engagement en lugar
de incluir todas las variables simultáneamente.

\subsubsection{Factor de inflación de varianza (VIF)}
\label{sssec:vif}

El análisis VIF sobre el conjunto de predictores candidatos
revela que las variables académicas presentan los valores más
altos (Parcial~1: VIF~$= 13{,}7$; AM~2: VIF~$= 10{,}4$;
AM~1: VIF~$= 9{,}4$), mientras que las de engagement son más
moderadas (módulos únicos P1: VIF~$= 7{,}2$;
tiempo hasta P1: VIF~$= 4{,}0$;
tiempo por visita P1: VIF~$= 3{,}9$). Ningún valor supera el
umbral de 15 que produce inestabilidad severa, pero la magnitud
de algunos VIF justifica el uso de regularización L1/L2 en los
modelos de regresión logística (ver Sección~\ref{sec:metodologia}).

% ------------------------------------------------------------
\subsection{Síntesis: Decisiones de Diseño del Dataset}
\label{subsec:sintesis_datos}
% ------------------------------------------------------------

La Tabla~\ref{tab:resumen_decisiones} sintetiza las decisiones
de preprocesamiento tomadas a partir del análisis exploratorio.

\begin{table}[htbp]
  \centering
  \caption{Resumen de decisiones de diseño del dataset con
           justificación empírica e impacto cuantificado.}
  \label{tab:resumen_decisiones}
  \small
  \begin{tabular}{p{3.5cm}p{5.0cm}p{4.8cm}}
    \toprule
    \textbf{Decisión} & \textbf{Justificación} & \textbf{Impacto} \\
    \midrule
    Excluir RETIRADO y SIN~NOTA
      & Etiqueta no observable; sesgo de retiro endógeno;
        heterogeneidad semántica
      & $N$: 764~$\to$~554 \\[3pt]
    Excluir estudiantes sin engagement
      & Cero ambiguo (75\,\% en semestre de arranque);
        $\Delta F_1 < 0{,}02$ al excluirlos
      & $N$: 554~$\to$~502 \\[3pt]
    Excluir variables de actividades evaluativas
      & AUC~$< 0{,}50$; $p > 0{,}05$ (sin discriminación)
      & $-3$ predictores candidatos \\[3pt]
    Incluir \texttt{era\_encoded}
      & $\Delta$AUC(P1) $= -0{,}087$;
        engagement mediano cae 60\,\% en era nueva
      & Control en todos los modelos \\[3pt]
    Incluir \texttt{modalidad\_am}
      & AM-3 grupal vs.\ individual:
        $p = 1{,}4\times10^{-9}$, $d = 0{,}40$
      & Control en modelos con AMs \\[3pt]
    Incluir engagement a pesar de efecto pequeño
      & $\rho_{\text{parcial}} = 0{,}166$, $p < 0{,}001$;
        valor incremental en zona marginal de P1
      & $+6$ predictores en CP1 \\[3pt]
    Aplicar balanceo SMOTE / SMOTEENN
      & 74,3\,\% APROBÓ vs.\ 25,7\,\% REPROBÓ;
        Recall sin balanceo $\approx 0{,}55$
      & Solo en training fold; nunca en evaluación \\
    \bottomrule
  \end{tabular}
\end{table}

El dataset de modelado resultante (\texttt{df\_activos\_engagement},
$N = 502$) presenta tres características estructurales que
condicionan el diseño metodológico de la siguiente sección:
(i)~\textbf{desbalance de clases} (74,3\,\% vs.\ 25,7\,\%),
que exige técnicas de remuestreo para evitar que los modelos
sesguen la predicción hacia la clase mayoritaria;
(ii)~\textbf{estructura longitudinal multisemestral}, que
invalida la validación cruzada aleatoria estándar y exige el
esquema LOSO; y (iii)~\textbf{heterogeneidad de era pedagógica},
que requiere variables de control y análisis estratificado.
Estos tres aspectos se abordan en detalle en la sección de
metodología.

% ============================================================
%  REFERENCIAS
% ============================================================

\begin{thebibliography}{9}

\bibitem{Lopez2021}
J.~López-Zambrano, J.~A. Lara~Torralbo, y C.~Romero,
``Early prediction of student learning performance through data
mining: A systematic review,''
\textit{Psicothema}, vol.~33, no.~3, pp.~456--465, 2021.

\bibitem{Tinto1993}
V.~Tinto,
\textit{Leaving College: Rethinking the Causes and Cures of
Student Attrition}, 2.ª~ed.
Chicago, IL, USA: University of Chicago Press, 1993.

\bibitem{Kerby2014}
D.~S. Kerby,
``The simple difference formula: An approach to teaching
nonparametric correlation,''
\textit{Compr. Psychol.}, vol.~3, p.~11.IT.3.4.1, 2014.

\bibitem{He2009}
H.~He y E.~A. Garcia,
``Learning from imbalanced data,''
\textit{IEEE Trans. Knowl. Data Eng.}, vol.~21, no.~9,
pp.~1263--1284, sep. 2009.

\end{thebibliography}

\end{document}
